% !TeX root = ./../main.tex

\chapter{Klasyfikacja w obecności szumu i wpływ redukcji błędów}\label{chapter:classification}

\agentinfo{
Rozdział bada, jak szum wpływa na klasyfikację prowadzoną z użyciem obwodów kwantowych i czy klasyczna redukcja błędów ją poprawia.
Spina tytuł pracy z rozdziałami 4--7.
Docelowa objętość: około 16 stron.
Decyzja autora z 2026-11-27: rozdział nie ma własnej tezy (robocza teza T4 odpada); odpowiada na pytanie badawcze o wpływ szumu i klasycznej korekty wyniku na klasyfikację, a odpowiedź jest warunkowa (sekcja 8.4.2). W sekcjach 8.1 i 8.5 nie odwoływać się do tezy, tylko do pytania badawczego; numer pytania uzgodnić po zatwierdzeniu wprowadzenia.
Struktura przyjęta przez autora 2026-11-26; eksperymenty E0--E4 są zadaniami w infokuźni (epika Eksperymenty uzupełniające).

Materiał istniejący: repozytorium /Users/krzysztofwerner/git/quantum-classifier (tylko do odczytu; nie czytać plików .env).
Klasyfikator FBS z odczytem klasy z fazy (azymutu wektora Blocha), kąty bramek przewiduje sieć klasyczna, 2--7 kubitów.
Dane istniejące są syntetyczne i tabelaryczne, nie obrazowe.
Są wyniki symulatorowe z szumem, przebiegi na urządzeniu ibm_fez, kalibracja modelu szumu stanami GHZ oraz afiniczny model urządzenia; nie zastosowano żadnej mitygacji błędów.
Zapisane pomiary sprzętowe (400 próbek uczących i 100 testowych na model) pozwalają badać korektę odczytu bez nowych przebiegów na sprzęcie.
Manuskrypt roboczy Quantum_Specific_Classification_working_draft_v4_kw.tex jest współautorski i nieopublikowany; część teoretyczna (geometria) jest wkładem współautora.

Nowe eksperymenty:
E0. Korekta odczytu typu PDU (interpolacja błędu fazy na kubit) dopasowana na próbkach uczących i oceniona na testowych, na zapisanych danych z ibm_fez.
E1. Krzywe degradacji dokładności przy rosnącym szumie na symulatorze, z korektą i bez niej, w tym skończona liczba pomiarów.
E2. Porównanie treningu w szumie z korektą po fakcie, na wielu ziarnach losowych.
E4. Klasyfikacja obrazów 8x8 lub 16x16 w wariantach: oryginał, rekonstrukcja LPIQE surowa, po PDU, po pix2pix, po obu etapach; klasyfikator klasyczny i klasyfikator FBS.
Oryginalnego kodu potoku LPIQE, PDU i pix2pix nie ma na komputerze autora; potok jest odtwarzany na symulatorze w ramach osobnego zadania.

Zagrożenia do usunięcia przed użyciem wyników: jedno ziarno losowe, brak zbioru walidacyjnego (dobór hiperparametrów i punktu kontrolnego na zbiorze testowym), standaryzacja liczona przed podziałem danych.
Klasyfikator FBS jest słabszy od sieci MLP we wszystkich przypadkach; w rozprawie jest środowiskiem do badania szumu i redukcji błędów, a nie propozycją konkurencyjnego klasyfikatora, i należy to napisać wprost.
Doprecyzowanie po sprawdzeniu plików (PHD-TS-78): zdanie o przewadze sieci MLP jest prawdziwe dla dokładności na pełnym zbiorze testowym w dziesięciu modelach M2--M11 (różnica od 0,011 do 0,081, czyli od 22 do 161 próbek z 2000; pliki experiment.json przebiegów testowych); na podzbiorach 100 próbek klasyfikator bywa lepszy (M9: 0,62 wobec 0,55; M11: 0,82 wobec 0,80), a manuskrypt podaje dwa wyjątki poza tymi modelami (komórka 5 klas, 300 cech, wariant Pauliego w zestawieniu symulatorowym oraz wariant z enkoderem transformerowym).
Rozstrzygnięcie na plikach (PHD-TS-79, opisane w 8.4.1): oba wyjątki potwierdzone (0,631 wobec 0,624 przy 37 z 50 zbiorów na korzyść klasyfikatora; 0,5555 wobec 0,5195 dla odmiany z siecią typu transformer); w zestawieniu symulatorowym klasyfikator ma niższą średnią w 31 z 32 konfiguracji i jest słabszy w 1239 z 1600 porównań na pojedynczych zbiorach, przy czym odniesieniem jest tam sieć MLP o jednej warstwie ukrytej z 64 neuronami, a nie sieć o warstwach klasyfikatora; sformułowanie we wszystkich przypadkach jest więc za mocne, właściwe jest: w porównaniach z siecią o tych samych warstwach na pełnym zbiorze testowym.

Stan (PHD-TS-88, 2026-11-27): sekcje 8.1--8.3 i podrozdziały 8.4.1--8.4.2 napisane, czekają na przegląd autora; podrozdział 8.4.3 napisany 2026-11-27 na prośbę autora z eksperymentu E2 (X5 i X5b), razem z akapitem metodyki w 8.3 i zdaniem o zakresie w 8.1; podrozdział 8.4.4 i sekcje 8.5--8.6 do napisania (czekają na eksperymenty).
Ustalenia nazewnicze dla sekcji 8.4--8.6:
Model: klasyfikator z odczytem fazowym, w skrócie klasyfikator FBS (nazwa z manuskryptu); klasyczny punkt odniesienia: sieć MLP o tych samych warstwach ukrytych.
Symbole: K liczba klas, d liczba cech, n liczba kubitów, L liczba bloków (jak w rozdziale 7; uwaga: w manuskrypcie i w kodzie jest odwrotnie, L oznacza tam liczbę klas, a K liczbę bloków), makro nshots liczba pomiarów na bazę; wektor Blocha j-tego kubitu to makro blochvec z indeksem j, jego azymut varphi_j, azymut obcięty varphi_j z kreską, rzut na płaszczyznę xy r_{j,xy}, promień to norma wektora; siły szumu na blok to makra cohangle (odstrojenie), qphase (tłumienie fazy), qdep (depolaryzacja), współczynnik skrócenia 1 - 4 qdep / 3 to makro distshrink.
Nie używać dla szumu liter epsilon, gamma i p z kodu i z plików README eksperymentów: epsilon to funkcja PDU, gamma to funkcja referencyjna.
Litera K oznacza liczbę klas, więc liczbę węzłów funkcji korekty (w README: K) oznacza w 8.4.2 makro nnodes; funkcję błędu fazy klasyfikatora oznacza makro phaseerr (litera nu, nie epsilon).
Warianty korekty odczytu (8.4.2): FO i FZ to funkcja błędu fazy odczytywana dla fazy oczekiwanej i dla fazy zmierzonej, AO to odwrócenie odwzorowania afinicznego, AR to regresja wektora odniesienia na wektor zmierzony (w plikach eksperymentów: a_pdu, a2_pdu_slepa, b_afin, b2_afin_wprost); wariant FO wprowadza do korekty informację o odpowiedzi i jest podawany tylko z wynikiem kontroli.
Modele: M1--M11 według tabeli tab:cls:models (kolejność malejącej liczby bramek CZ), z identyfikatorami katalogów eksperymentów o budowie c7_f200_q3_b13 (klasy, cechy, kubity, bloki); odpowiedniki w manuskrypcie: model A = M1, model B = M2, model C = M6; M2 jako jedyny odczytuje dwie współrzędne z kubitu (przyrostek _d2).
Warianty bloku: wariant Pauliego (w kodzie pauli) i wariant uniwersalny (w kodzie universal).
Zbiory: typ problemu to para (K, d); seria trudności to 50 zbiorów jednego typu (katalog Dif_compl_data); modele M1--M11 uczono na pierwszym zbiorze serii (dataset_00); w danych jest też seria 50 zbiorów losowych (Rand_data), użyta tylko w 2 z 58 przebiegów wsadowych i w rozdziale pominięta.
Poziomy szumu: poziomy stopniowane (0; 0,01; 0,05; 0,10), poziom ibm_fez (trójka z kalibracji stanami GHZ, kolumny tabeli tab:cls:models), krotność poziomu ibm_fez (przebieg mieszany), siatka sił eksperymentu symulacyjnego (0,005--0,5).
Przebiegi na urządzeniu: przebieg testowy (100 próbek testowych), przebieg kalibracyjny (400 próbek uczących, pierwsza runda dostrajania, model niezmieniony), dostrajanie na urządzeniu (trzy rundy).
Eksperymenty: E0 to w tekście eksperyment na zapisanych pomiarach, E1 to eksperyment symulacyjny, E2 to eksperyment porównawczy (konfiguracja płytka = konfiguracja przeglądu szumu, konfiguracja głęboka = model M2 uczony od nowa); E4 nie ma jeszcze nazwy w tekście.
Miary: dokładność, zgodność predykcji, odległość wektorów Blocha, średni promień, błąd azymutu, odzysk; wartości dokładne oznaczają granicę nieskończonej liczby pomiarów.
Terminologia: korekta odczytu (nie korekcja); zagrożenia trafności mają etykiety itm:cls:ZT1 do itm:cls:ZT5 (ZT1 wybór modelu i standaryzacja z udziałem zbioru testowego, ZT2 jedno ziarno i jeden zbiór, ZT3 mała próba z urządzenia, ZT4 model szumu, ZT5 kalibracja korekty na próbkach uczących i względem symulatora).
Etykiety dodane w sekcjach 8.1--8.3: fig:cls:model (panele fig:cls:circuit i fig:cls:readout), eq:cls:readout, eq:cls:azimuth-noise, eq:cls:noise-frame, tab:cls:models, itm:cls:ZT1 do itm:cls:ZT5.
Etykiety dodane w 8.4.1: tab:cls:hw-results (modele M1--M11: symulator, urządzenie, dostrajanie), fig:cls:noise-results (panel a: promień i zgodność predykcji a liczba bramek CZ; panel b: przegląd szumu); nazwa wprowadzona w 8.4.1: przegląd szumu (64 przebiegi uczenia jednej konfiguracji przy kombinacjach poziomów stopniowanych).
Objętość sekcji 8.1--8.3 w kompilacji bez notatek: około 6,9 strony (8.1: 0,9; 8.2: 2,4; 8.3: 3,6; planowano łącznie około 5); kandydaci do cięcia, jeśli budżet rozdziału tego wymaga: tabela tab:cls:models do dodatku (0,4 strony), opis kalibracji GHZ, lista zagrożeń trafności zapisana jako jeden akapit, szczegóły siatki eksperymentu symulacyjnego przeniesione do sekcji 8.4.2.
Objętość podrozdziału 8.4.1 w kompilacji bez notatek: około 3,7 strony (tekst 2,9; tabela tab:cls:hw-results 0,4; rysunek 0,35; planowano około 3); kandydaci do cięcia: tabela do dodatku (0,4 strony), zdania o symulacji na fałszywym urządzeniu i o współczynniku determinacji odwzorowań (0,15), zdanie o stanach przygotowanych jedną warstwą obrotów (0,1), opis rundy dostrajania (0,15).
Etykiety dodane w 8.4.2: eq:cls:corrections (cztery warianty korekty), tab:cls:pdu-recovery (odzysk na progu spadku dokładności, modele M1--M3), fig:cls:pdu-curves (krzywe modelu M2).
Objętość podrozdziału 8.4.2 w kompilacji bez notatek: około 3,7 strony (tekst 3,1; tabela 0,27; rysunek 0,33; planowano około 3); kandydaci do cięcia są w konspekcie podrozdziału.
Odwołania zapisywać jak w pozostałych rozdziałach, słowem i poleceniem ref (cleveref nie ma polskich nazw i nie jest w rozprawie używany).
}

\section{Cel i zakres rozdziału}\label{sec:cls:scope}
\agentinfo{
Pytanie rozdziału sformułowano opisowo, bez numeru i bez brzmienia tezy: robocze oznaczenia PB5 i T4 z konspektu wprowadzenia nie są zatwierdzone przez autora.
Wynik może być negatywny i wtedy tezę trzeba przeformułować, a nie naginać wyniki.
Zakres obejmuje tylko materiał istniejący i eksperymenty wykonane (E0, E1, E2); eksperymentu E4 tekst nie opisuje.
Podział wkładu oparto na oznaczeniach w manuskrypcie (fragmenty z komentarzem Filled in by K. Werner: sekcje 3.1, 3.5--3.7, 4.2--4.6 oraz część 3.2 o sieci MLP); manuskrypt nie ma listy autorów, więc nazwisk nie podano.
}

Rozdziały~\ref{chapter:pdu}--\ref{chapter:learned_encoding} dotyczyły obliczeń, których wynik jest wielkością ciągłą, a miarą jakości był błąd tego wyniku.
W klasyfikacji wynik obliczenia kwantowego jest krokiem pośrednim, a ocenie podlega decyzja, czyli klasa przypisana próbce.
Rozdział odpowiada na pytanie, jak szum obwodu i skończona liczba pomiarów wpływają na skuteczność klasyfikacji prowadzonej z użyciem obwodu kwantowego oraz czy klasyczna korekta wyniku, kalibrowana eksperymentalnie na wzór funkcji PDU (rozdział~\ref{chapter:pdu}), tę skuteczność przywraca; jest to jedno z pytań badawczych postawionych w podrozdziale~\ref{sec:introduction:theses}.
Odpowiedź nie jest przesądzona: decyzja jest wielkością dyskretną, więc błąd wyniku, który nie przenosi próbki przez granicę decyzyjną, nie obniża dokładności, a jego usunięcie niczego wtedy nie poprawia.
Rozdział dopuszcza zatem wynik negatywny i raportuje go na równi z pozytywnym.
\agenttask{Po zatwierdzeniu przez autora numeracji i brzmienia pytań badawczych oraz tez (robocze PB5 i T4 w konspekcie chapters/introduction.tex) uzgodnić z wprowadzeniem akapit powyżej: podać numer pytania i odwołać się do właściwej tezy. Robocze brzmienie PB5 mówi o klasyfikacji obrazów, a istniejące wyniki dotyczą danych tabelarycznych; rozbieżność zniknie dopiero po eksperymencie E4 albo po zmianie brzmienia pytania.}

Badania prowadzi się na hybrydowym klasyfikatorze z odczytem fazowym (podrozdział~\ref{sec:cls:classifier}), w którym klasę odczytuje się z faz lokalnych kubitów, a więc z wielkości, której zniekształcenia są przedmiotem metody PDU.
Klasyfikator nie jest propozycją metody konkurencyjnej wobec klasyfikatorów klasycznych: we wszystkich dziesięciu modelach, dla których wyznaczono klasyczny punkt odniesienia, jego dokładność na zbiorze testowym jest o $0{,}011$ do $0{,}081$ niższa od dokładności sieci MLP o tych samych warstwach ukrytych.
% ŹRÓDŁO: /Users/krzysztofwerner/git/quantum-classifier/runs/*/hardware_ibm_fez_*/experiment.json modeli M2--M11, pola checkpoint.test_accuracy i mlp.best_test_accuracy (pełny zbiór testowy, 2000 próbek); skrajne różnice: M11 0,8065 wobec 0,8175 (22 próbki, 0,011) oraz M2 0,4390 wobec 0,5195 (161 próbek, 0,0805, po zaokrągleniu 0,081).
Służy wyłącznie jako układ, w którym wpływ szumu na decyzję i działanie korekty można zmierzyć na symulatorze i na rzeczywistym urządzeniu.

Zakres rozdziału obejmuje istniejące wyniki klasyfikatora na symulatorze z modelem szumu, zapisane wyniki przebiegów na urządzeniu ibm\_fez oraz trzy eksperymenty wykonane na potrzeby rozprawy: korektę odczytu dopasowaną i ocenioną na zapisanych pomiarach z urządzenia, krzywe degradacji dokładności przy rosnącym szumie na symulatorze, z korektą i bez niej, oraz porównanie korekty z treningiem w szumie na pięciu ziarnach losowych.
Dane są syntetyczne i tabelaryczne, a nowych przebiegów na urządzeniu nie wykonywano.
\agenttask{Po wykonaniu eksperymentu E4 (klasyfikacja obrazów po redukcji błędów; nierozpoczęty) dopisać go do akapitu o zakresie i uzupełnić zdanie o danych. Do tego czasu nie opisywać go jako wykonanego. Eksperyment E2 dopisano 2026-11-27.}

Klasyfikator opisano w nieopublikowanym manuskrypcie roboczym, przygotowywanym przez autora rozprawy wspólnie z innymi osobami: geometryczne sformułowanie zadania klasyfikacji jest wkładem współautora, zbiory danych dostarczyli współpracownicy, a udziałem autora rozprawy są implementacja modelu, model szumu i jego kalibracja oraz eksperymenty na symulatorze i na urządzeniu.
Wkładem tego rozdziału jest analiza klasyfikatora od strony redukcji błędów: zestawienie odczytu fazowego z założeniami metody PDU, eksperymenty z korektą odczytu, w tym jej porównanie z treningiem w szumie, i jawne wskazanie zagrożeń trafności istniejących wyników.
\agenttask{Pytania do autora w sprawie autorstwa (źródła tego nie rozstrzygają): 1. Kto jest autorem części teoretycznej manuskryptu (rozmaitość stanów klasyfikacji, reprezentacja Delta, wiązka reperów) i jak go wskazać w tekście? 2. Czyim pomysłem są odczyt klasy z azymutu wektora Blocha, promień jako miara pewności i budowa bloku ansatzu (dokument docs/design-proposal.md w repozytorium klasyfikatora przypisuje je projektowi implementacji)? 3. Jak cytować manuskrypt: czy dodać wpis typu unpublished do bibliography/own_publications.bib (potrzebne tytuł i lista autorów)? 4. Kto wygenerował zbiory danych i jakim narzędziem?}

\section{Klasyfikator z odczytem fazowym}\label{sec:cls:classifier}
\agentinfo{
Opis zwięzły, w zakresie potrzebnym do zrozumienia wpływu szumu; geometria modelu przypisana współautorowi manuskryptu (nazwisko do uzupełnienia przez autora).
Źródła: kod \agentref{/Users/krzysztofwerner/git/quantum-classifier/torus_classifier/circuit.py} oraz moduły readout.py, simplex.py, losses.py i encoder.py z tego samego katalogu; manuskrypt roboczy w katalogu \agentref{/Users/krzysztofwerner/git/quantum-classifier}, plik Quantum_Specific_Classification_working_draft_v4_kw.tex (sekcje 2.3, 2.4, 2.7 i 3.6).
Liczby kątów na blok pochodzą z kodu (circuit.n_block_params); dla n = 6 jest to 78 (wariant Pauliego) i 198 (wariant uniwersalny); wzór M = 3KQ^2 z sekcji 3.6 manuskryptu (część niepisana przez autora rozprawy) nie zgadza się z kodem i nie jest używany.
Zapis odczytu: w kodzie i w manuskrypcie współrzędna nazywa się delta i jest azymutem podzielonym przez pi; w rozprawie wzory zapisano wprost przez azymut, bo litera delta oznacza w rozdziale 4 dryf błędu.
Azymut obcina się w implementacji z marginesem 0,001 pi od krańców przedziału; sieć klasyczna ma funkcję aktywacji ReLU; składnik fazowy straty liczy się ze składowych wektora, bez funkcji arg (te szczegóły pominięto w tekście dla budżetu stron).
Dodane w rozprawie, nieobecne w źródłach (do akceptacji autora): oszacowanie niepewności azymutu eq:cls:azimuth-noise oraz zestawienie odczytu fazowego z założeniami Z3 i Z7 metody PDU.
Definicje wariantów korekty odczytu należą do sekcji 8.4.2; tutaj jest tylko uzasadnienie, dlaczego warianty są potrzebne.
}

Klasyfikator pochodzi z nieopublikowanego manuskryptu roboczego, w którym nosi nazwę FBS (ang. \english{Frame Bundle Section}).
Jego punktem wyjścia jest geometryczne sformułowanie zadania klasyfikacji, będące wkładem współautora manuskryptu: rozkład prawdopodobieństwa $K$ klas ma $K-1$ niezależnych współrzędnych, każdą z nich zamienia się na kąt, a klasyfikator opisuje się lokalnie jako cięcie wiązki reperów nad rozmaitością stanów klasyfikacji.
Rozprawa korzysta tylko z obliczeniowej konsekwencji tego sformułowania: $K-1$ kątów można zapisać w fazach lokalnych $K-1$ kubitów i odczytać z kierunków ich wektorów Blocha.

\paragraph{Sieć klasyczna i obwód.}
Model jest \glslink{hybrid-model}{modelem hybrydowym} z obwodem o stałej strukturze (podrozdział~\ref{subsec:tb:ml:qml}, rysunek~\ref{fig:tb:ml:hybrid}): sieć klasyczna $\Theta_{\mathbf{w}}$, perceptron wielowarstwowy z wyjściem $\pi\tanh(\cdot)$, zamienia wektor $d$ cech próbki na kąty wszystkich bramek, $\boldsymbol{\theta}=\Theta_{\mathbf{w}}(\mathbf{x})$.
Obwód nie zawiera bramek kodujących dane ani własnych parametrów uczonych, więc próbka wpływa na stan kwantowy wyłącznie przez kąty.
Na $n$ kubitów w stanie $\ket{0}$ działają kolejno warstwa bramek Hadamarda, $L$ bloków o jednakowej budowie i druga warstwa bramek Hadamarda (rysunek~\ref{fig:cls:circuit}).
Blok składa się z bramki jednokubitowej o trzech kątach na każdym kubicie oraz, dla każdej uporządkowanej pary kubitów, dwóch bramek dwukubitowych sterowanych w przeciwnych kierunkach, razem $2n(n-1)$ bramek dwukubitowych.
W wariancie Pauliego bramkami są złożenie obrotów $R_x$, $R_y$, $R_z$ i kontrolowane przesunięcie fazy $\mathrm{CP}$, co daje $n(2n+1)$ kątów na blok, a w wariancie uniwersalnym ogólny obrót jednokubitowy (bramka $\mathrm{U3}$ z dokładnością do fazy globalnej) i jego wersja kontrolowana, czyli $3n(2n-1)$ kątów; wszystkie modele uruchomione na urządzeniu (tabela~\ref{tab:cls:models}) używają wariantu uniwersalnego.
% ŹRÓDŁO: /Users/krzysztofwerner/git/quantum-classifier/torus_classifier/circuit.py (funkcje ansatz_block i n_block_params) oraz encoder.py (wyjście pi * tanh).

\begin{figure}[tb]
\centering
\begin{subfigure}[c]{0.47\textwidth}
\centering
% !TeX root = ./../../main.tex
% Obwód klasyfikatora z odczytem fazowym dla n = 2 kubitów (K = 3 klasy); rys. fig:cls:circuit (panel a rysunku fig:cls:model), podrozdział 8.2.
% Rysunek własny (quantikz2, zadanie PHD-TS-78), wykonany na podstawie kodu
% /Users/krzysztofwerner/git/quantum-classifier/torus_classifier/circuit.py
% (funkcje _apply_sandwiched_ansatz, ansatz_block, _apply_block_noise).
% Kolejność bramek dwukubitowych w bloku jak w ansatz_block: dla każdej uporządkowanej pary (j, m), j != m,
% bramka sterowana kubitem j, a po niej bramka sterowana kubitem m; dla n = 2 daje to cztery bramki.
% Kanał błędu (linia kropkowana) występuje tylko w symulacji z modelem szumu.
% Każdy kubit jest mierzony w trzech bazach (sigma_x, sigma_y, sigma_z), co daje jego wektor Blocha.
\begin{quantikz}[row sep=0.35cm, column sep=0.13cm, font=\footnotesize]
\lstick{$\ket{0}$} & \gate{\mathrm{H}} & \gate{U}\gategroup[2,steps=6,style={dashed,rounded corners,inner xsep=1pt,inner ysep=3pt},label style={yshift=1pt}]{blok $\ell=1,\dots,L$} & \ctrl{1} & \gate{V} & \gate{V} & \ctrl{1} & \gate[style={densely dotted}]{\errch} & \gate{\mathrm{H}} & \meter{} \\
\lstick{$\ket{0}$} & \gate{\mathrm{H}} & \gate{U} & \gate{V} & \ctrl{-1} & \ctrl{-1} & \gate{V} & \gate[style={densely dotted}]{\errch} & \gate{\mathrm{H}} & \meter{}
\end{quantikz}

\caption{}\label{fig:cls:circuit}
\end{subfigure}\hfill
\begin{subfigure}[c]{0.52\textwidth}
\centering
% !TeX root = ./../../main.tex
% Odczyt klasy z kierunku wektora Blocha: rzut na płaszczyznę xy sfery Blocha kubitu j;
% rys. fig:cls:readout (panel b rysunku fig:cls:model), podrozdział 8.2.
% Rysunek własny (TikZ, zadanie PHD-TS-78), wykonany na podstawie kodu
% /Users/krzysztofwerner/git/quantum-classifier/torus_classifier/readout.py (azymut, obcięcie, promień)
% i torus_classifier/losses.py (kierunki docelowe w azymucie -pi/2 i +pi/2).
% Wektory przykładowe: azymut -50 stopni, długość rzutu 0,85 (bez szumu, cienka strzałka) i 0,36 (po skróceniu, gruba szara strzałka);
% wektor po zniekształceniu fazy (strzałka przerywana): azymut +25 stopni, długość 0,85.
% Kółka wokół końców wektorów mają jednakowy promień (niepewność statystyczna składowych, schematycznie, nie w skali).
\begin{tikzpicture}[
    scale=1.7,
    >={Stealth[length=2mm]},
    aux/.style={thin, gray},
    lab/.style={font=\footnotesize, align=left, inner sep=1.5pt},
  ]
  % półpłaszczyzna r_x < 0: azymut obcinany do +-pi/2
  \begin{scope}
    \clip (0,0) circle[radius=1];
    \fill[gray!22] (-1.05,-1.05) rectangle (0,1.05);
  \end{scope}
  \draw[thin] (0,0) circle[radius=1];
  \node[lab, align=center, font=\scriptsize] at (-0.47,0.27) {nasycenie\\$r_{j,x}<0$};
  % osie
  \draw[->, thin] (-1.15,0) -- (1.45,0) node[above, font=\footnotesize] {$r_{j,x}$};
  \draw[->, thin] (0,-1.25) -- (0,1.38) node[left, font=\footnotesize] {$r_{j,y}$};
  % granica decyzji dla K = 2
  \draw[very thick, densely dashed] (0,0) -- (1,0);
  % kierunki docelowe
  \fill (0,-1) circle[radius=0.9pt];
  \fill (0,1) circle[radius=0.9pt];
  \node[lab, anchor=south west] at (0.06,1.03) {$\varphi_j=+\tfrac{\pi}{2}$: nie klasa $j$};
  \node[lab, anchor=north west] at (0.06,-1.03) {$\varphi_j=-\tfrac{\pi}{2}$: klasa $j$};
  % punkty końcowe wektorów
  \coordinate (A) at ({0.85*cos(-50)},{0.85*sin(-50)});
  \coordinate (S) at ({0.36*cos(-50)},{0.36*sin(-50)});
  \coordinate (R) at ({0.85*cos(25)},{0.85*sin(25)});
  % wektor bez szumu: cienka strzałka
  \draw[->, line width=0.6pt] (0,0) -- (A);
  % wektor skrócony (szum niekoherentny): ten sam azymut, gruba krótka strzałka na wierzchu
  \draw[-{Stealth[length=2.8mm,width=2.4mm]}, line width=2pt, black!55] (0,0) -- (S);
  % niepewność statystyczna składowych (schematycznie)
  \draw[thin, densely dotted] (A) circle[radius=0.095];
  \draw[thin, densely dotted] (S) circle[radius=0.095];
  % wektor po zniekształceniu fazy: inny azymut
  \draw[->, line width=0.9pt, densely dashed] (0,0) -- (R);
  % kąty
  \draw[thin] (0.2,0) arc[start angle=0, end angle=-50, radius=0.2];
  \node[font=\footnotesize] at ({0.34*cos(-19)},{0.34*sin(-19)}) {$\varphi_j$};
  \draw[->, thin] ({0.62*cos(-50)},{0.62*sin(-50)}) arc[start angle=-50, end angle=25, radius=0.62];
  % opisy wektorów
  \node[lab, anchor=west] (lr) at (1.3,0.72) {zniekształcenie fazy:\\azymut przesunięty};
  \draw[aux] (lr.west) -- ($(R)+(0.04,0.03)$);
  \node[lab, anchor=west] (ls) at (1.3,-0.36) {skrócenie wektora:\\azymut bez zmian};
  \draw[aux] (ls.west) -- ($(S)+(0.1,-0.02)$);
  \node[lab, anchor=west] (la) at (1.3,-0.8) {bez szumu};
  \draw[aux] (la.west) -- ($(A)+(0.1,-0.03)$);
\end{tikzpicture}

\caption{}\label{fig:cls:readout}
\end{subfigure}
\caption[Klasyfikator z odczytem fazowym: obwód i odczyt klasy]{Klasyfikator z odczytem fazowym. (a)~Obwód dla $n=2$ kubitów ($K=3$ klasy): $U$ oznacza bramkę jednokubitową, $V$ bramkę dwukubitową sterowaną zaznaczonym kubitem, a kanał błędu $\errch$ występuje tylko w symulacji z modelem szumu. (b)~Odczyt klasy z kierunku wektora Blocha $j$-tego kubitu w rzucie na płaszczyznę $xy$: punkty na osi $r_{j,y}$ są kierunkami wzorcowymi, gruba linia przerywana ($\varphi_j=0$) granicą decyzji dla $K=2$, w obszarze zacienionym azymut jest obcinany, a kółka oznaczają schematycznie niepewność statystyczną składowych. Opracowanie własne.}
\label{fig:cls:model}
\end{figure}

\paragraph{Odczyt klasy.}
Dla każdego kubitu wyznacza się wartości oczekiwane trzech macierzy Pauliego, czyli wektor Blocha $\blochvec_j=(r_{j,x},r_{j,y},r_{j,z})$ jego stanu zredukowanego (równanie~\eqref{eq:tb01:blochrho}), a z niego azymut $\varphi_j=\arg(r_{j,x}+\mathrm{i}\,r_{j,y})$, długość rzutu na płaszczyznę $xy$, $r_{j,xy}=(r_{j,x}^2+r_{j,y}^2)^{1/2}$, i promień $\|\blochvec_j\|$.
W odczycie podstawowym $n=K-1$.
Azymuty obcina się do przedziału $(-\tfrac{\pi}{2},\tfrac{\pi}{2})$, co daje wartości $\bar{\varphi}_j$, i zamienia na rozkład prawdopodobieństwa klas:
\begin{equation}\label{eq:cls:readout}
\hat{y}_k=\Bigl(\frac{1}{2}-\frac{\bar{\varphi}_k}{\pi}\Bigr)\prod_{i=1}^{k-1}\Bigl(\frac{1}{2}+\frac{\bar{\varphi}_i}{\pi}\Bigr)\quad\text{dla } k<K,\qquad
\hat{y}_K=\prod_{i=1}^{K-1}\Bigl(\frac{1}{2}+\frac{\bar{\varphi}_i}{\pi}\Bigr),
\end{equation}
a predykcją jest klasa o największej wartości $\hat{y}_k$.
Czynnik $\tfrac{1}{2}-\bar{\varphi}_j/\pi$ jest prawdopodobieństwem warunkowym przynależności do klasy $j$, jeśli próbka nie należy do klas o niższych numerach, więc wzorcem uczącym klasy $k<K$ są azymuty $\varphi_k=-\tfrac{\pi}{2}$ i $\varphi_i=+\tfrac{\pi}{2}$ dla $i\neq k$, a klasy $K$ wszystkie azymuty równe $+\tfrac{\pi}{2}$ (rysunek~\ref{fig:cls:readout}); dla $K=2$ granicą decyzji jest $\varphi_1=0$.
W modelu M2 z każdego kubitu odczytuje się dwie współrzędne: rolę kolejnych azymutów we wzorze~\eqref{eq:cls:readout} pełnią wielkości $2\theta_j-\pi$, gdzie $\theta_j$ jest kątem biegunowym wektora $\blochvec_j$, i siedem klas mieści się na trzech kubitach.
% ŹRÓDŁO: torus_classifier/readout.py (expectations_to_delta, polar_delta, bloch_to_readout), simplex.py (delta_to_prob), losses.py (labels_to_delta).
Strata próbki jest sumą składnika fazowego $\sum_j\bigl[1-\cos(\varphi_j-\varphi_j^{\ast})\bigr]$, w którym $\varphi_j^{\ast}$ jest azymutem wzorcowym, składnika $\sum_j(1-\|\blochvec_j\|)^2$ z wagą $0{,}01$ i entropii krzyżowej rozkładu~\eqref{eq:cls:readout} z wagą $0{,}5$, a uczenie prowadzi się na symulatorze PennyLane~\cite{bergholm_pennylane_2018}, na dokładnych wartościach oczekiwanych (granica $\nshots\to\infty$).
% ŹRÓDŁO: torus_classifier/losses.py (CompositeTorusLoss) i config.py (lambda_conf = 0.01, lambda_ce = 0.5); wartości potwierdzone w plikach experiment.json modeli M1--M11.

\paragraph{Promień jako miara pewności.}
Rozkład~\eqref{eq:cls:readout} zależy wyłącznie od kierunków wektorów Blocha, a promień, mniejszy od jedności, gdy kubit jest splątany z pozostałymi albo gdy stan został zaszumiony (podrozdział~\ref{subsec:tb:density}), określa jedynie, jak pewnie można decyzję odczytać.
Każdą składową wektora estymuje się z $\nshots$ pomiarów z odchyleniem standardowym nie większym niż $1/\sqrt{\nshots}$ (równanie~\eqref{eq:tb01:estimator}), więc w przybliżeniu liniowym estymator azymutu $\tilde{\varphi}_j$ ma odchylenie standardowe
\begin{equation}\label{eq:cls:azimuth-noise}
\sqrt{\mathrm{Var}\,\tilde{\varphi}_j}\;\lesssim\;\frac{1}{r_{j,xy}\sqrt{\nshots}},
\end{equation}
czyli około $0{,}31$~rad ($18^{\circ}$) dla $\nshots=1024$ i $r_{j,xy}=0{,}1$.
% ŹRÓDŁO: obliczenie własne: Var(phi) = (x^2 Var(y) + y^2 Var(x)) / r_xy^4 <= 1/(N r_xy^2) dla niezależnych estymatorów składowych; 1/(0,1 * 32) = 0,3125 rad = 17,9 stopnia.
Szum, który tylko skraca wektor, zwiększa więc rozrzut azymutu, a szum zmieniający kierunek wektora przesuwa azymut systematycznie (rysunek~\ref{fig:cls:readout}).

\paragraph{Związek z metodą PDU.}
W obwodzie referencyjnym metody PDU argument funkcji jest zapisany w fazie lokalnej kubitu odczytu, a wynikiem jest prawdopodobieństwo, w które tę fazę zamienia interferencja (podrozdział~\ref{sec:pdu:reference}).
W klasyfikatorze wynikiem jest sama faza: jej zniekształcenie przesuwa wprost współrzędną, od której zależy decyzja, podczas gdy utrata czystości stanu zmienia tylko promień.
Klasyfikator jest dzięki temu naturalnym układem do sprawdzenia, czy korekta wyniku kalibrowana eksperymentalnie przenosi się na zadanie decyzyjne, ale odczyt fazowy nie spełnia dwóch założeń metody (podrozdział~\ref{sec:pdu:assumptions}).
Po pierwsze, wartości oczekiwane nie są znane niezależnie od obwodu (założenie~\ref{itm:pdu:Z3}): odniesieniem może być tylko klasyczna symulacja tego samego obwodu, wykonalna dla badanych rejestrów o co najwyżej siedmiu kubitach, a wyznaczony błąd jest odchyleniem od symulatora, nie od wartości idealnej.
Po drugie, nie jest znany argument korekty (założenie~\ref{itm:pdu:Z7}): faza kubitu zależy od danych i jest wielkością, którą odczyt ma dopiero wyznaczyć, co odpowiada przypadkowi wyłączonemu z rozważań w podrozdziale~\ref{par:pdu:known-argument}.
Funkcję błędu można więc odczytywać albo dla fazy oczekiwanej z symulacji ocenianej próbki, co wprowadza do korekty informację o poprawnej odpowiedzi, albo dla fazy zmierzonej, obarczonej błędem, który ma zostać usunięty, a kalibrować ją trzeba nie na zaprojektowanej siatce argumentów, lecz na próbkach danych, których fazy skupiają się wokół kierunków wzorcowych.
Warianty korekty odczytu wynikające z tych wyborów definiuje podrozdział~\ref{subsec:cls:results:pdu}.

\section{Metodyka badań eksperymentalnych}\label{sec:cls:methodology}
\agentinfo{
Sekcja opisuje metodykę wyników istniejących oraz eksperymentów wykonanych E0, E1 i E2; metodyka poprawiona (zbiór walidacyjny, standaryzacja po podziale, pięć ziaren) dotyczy tylko eksperymentu E2 (akapit Eksperyment porównawczy, dopisany 2026-11-27).
Dane obrazowe (eksperyment E4) nie są jeszcze opisane.
Wszystkie parametry sprawdzono w kodzie i w plikach JSON przebiegów, nie tylko w manuskrypcie: \agentref{/Users/krzysztofwerner/git/quantum-classifier/torus_classifier/data.py} oraz config.py, hardware.py, hardware_training.py i ibm_fez_ghz_calibrations.json z tego samego katalogu, pliki hyperparameters.json i experiment.json przebiegów modeli M1--M11, pliki case_info.txt i complexity.csv w katalogu data.
Metodyka eksperymentów E0 i E1: \agentref{eksperymenty/pdu-ibm-fez/klasyfikator/README.md} \agentref{eksperymenty/krzywe-degradacji/klasyfikator/README.md}.
Rozbieżności manuskrypt--kod (w tekście użyto wartości z kodu i danych): manuskrypt podaje 1800 zbiorów w 18 typach (2--7 klas, 100--300 cech, po 100 zbiorów) i podział 80:20, a w danych i kodzie są 32 typy (3--10 klas, 100--400 cech) po 50 zbiorów serii trudności i 50 zbiorów losowych oraz podział 60:40; manuskrypt mówi o zbiorze walidacyjnym, którego w kodzie nie ma; manuskrypt mówi o tomografii stanu GHZ, a wierność wyznaczono metodą koherencji wielokwantowych.
Pominięte w tekście dla budżetu stron: problem dwuklasowy otrzymany z trójklasowego przez usunięcie jednej klasy (katalog Case_F100_C2_from_C3_pair01); zakresy dopasowanych sił szumu (depolaryzacja 0,0030--0,0053, odstrojenie 1,3--9,0 mrad, tłumienie fazy 0 z wyjątkiem 0,00015 dla trzech kubitów); daty kalibracji GHZ (10 września 2026 dla sześciu kubitów, 5 listopada 2026 dla pozostałych); szczegóły symulatorów odniesienia w przebiegu testowym (wartości dokładne i 1024 pomiary).
Procedurę kalibracji GHZ opisano na podstawie części manuskryptu napisanej przez autora rozprawy (sekcja 3.5) i komentarzy w config.py; kodu projektu GHZ (katalog poza repozytorium klasyfikatora) nie czytano.
Do opisania w sekcji 8.4.1 (tu pominięte dla budżetu stron): parametry uczenia (współczynnik uczenia 0,001 zmniejszany o połowę po pięciu epokach bez poprawy straty testowej, wsady po 32 próbki, ziarno 42) oraz przebieg rundy dostrajania na urządzeniu (pomiar 400 próbek uczących losowanych warstwowo z części uczącej, kontrola akceptacji na stracie zmierzonej, odwzorowanie afiniczne wektora Blocha symulowanego w zmierzony dopasowane na kubit metodą najmniejszych kwadratów, pięć epok optymalizatorem Adam ze współczynnikiem 0,0001, ponowny pomiar 100 próbek testowych); źródło: torus_classifier/hardware_training.py.
Wzór eq:cls:noise-frame wyprowadzono w rozprawie i sprawdzono rachunkiem na macierzach 2x2 (zgodność do 1e-15); README eksperymentu E1 podaje tę samą własność opisowo.
Podrozdział 2.2 (akapit o modelu warstwowym) podaje, że model używany w rozdziałach 7 i 8 zawiera błąd odczytu; w kodzie klasyfikatora błędu odczytu nie ma, co zaznaczono w tekście.
}

\paragraph{Dane i uczenie.}
Klasyfikator uczono na syntetycznych zbiorach tabelarycznych dostarczonych przez współpracowników.
Każdy zbiór ma $5000$ próbek, a typ problemu określają liczba cech $d\in\{100,200,300,400\}$ i liczba klas $K$ od $3$ do $10$; według metadanych $35$ cech jest informatywnych i $25$ redundantnych, a każdą klasę tworzą dwa skupienia.
Dla każdego typu dostępna jest seria $50$ zbiorów o stopniowanej trudności, w której separacja klas rośnie liniowo od $0{,}6$ do $1{,}1$, a udział losowo zmienionych etykiet maleje od $0{,}08$ do $0{,}01$; seria jest gradientem trudności, a nie zbiorem niezależnych powtórzeń.
Istniejące wyniki obejmują zestawienie symulatorowe, w którym jedną konfigurację modelu uczono kolejno na wszystkich zbiorach serii, oraz jedenaście modeli M1--M11 (tabela~\ref{tab:cls:models}), uczonych na pierwszym, najtrudniejszym zbiorze serii i uruchomionych na urządzeniu.
% ŹRÓDŁO: /Users/krzysztofwerner/git/quantum-classifier/data/Case_F*_C*/case_info.txt i complexity.csv (jednakowe dla wszystkich typów); pole data w plikach hyperparameters.json modeli M1--M11 (Dif_compl_data/dataset_00.npz); pliki summary.json przebiegów wsadowych runs/batch_* (n_datasets = 50, 56 z 58 przebiegów na serii Dif_compl_data).
\agenttask{Uzupełnić opis generatora danych po uzyskaniu informacji od autora albo współpracowników: nazwy parametrów w metadanych (class_sep, flip_y, cechy informatywne i redundantne, skupienia na klasę) odpowiadają funkcji make_classification biblioteki scikit-learn, ale kodu generatora nie ma w repozytorium, więc narzędzia nie nazwano i nie cytowano. Potrzebne także liczności klas.}
Cechy standaryzuje się średnią i odchyleniem standardowym policzonymi na całym zbiorze, po czym zbiór dzieli się losowo, z ziarnem $42$, na część uczącą ($3000$ próbek) i testową ($2000$ próbek); zbioru walidacyjnego nie ma.
% ŹRÓDŁO: torus_classifier/data.py, funkcja load_dataset (splits=(0.6, 0.4), seed=42); n_train i n_test w plikach hyperparameters.json.
Modele uczono optymalizatorem Adam (równanie~\eqref{eq:tb:ml:adam}) z jednym ziarnem generatora liczb losowych, a harmonogram współczynnika uczenia, punkt kontrolny i hiperparametry, czyli budowę sieci $\Theta_{\mathbf{w}}$ i w części modeli liczbę bloków, wybierano według wyniku na zbiorze testowym.
Klasycznym punktem odniesienia modeli M2--M11 jest sieć MLP o tych samych warstwach ukrytych i wyjściu softmax, uczona i wybierana tak samo; w zestawieniu symulatorowym odniesieniem jest sieć MLP o jednej warstwie ukrytej z $64$ neuronami, a model M1 nie ma własnego odniesienia.
% ŹRÓDŁO: pliki hyperparameters.json modeli M1--M11 (optimizer Adam, learning_rate 0.001, batch_size 32, scheduler ReduceLROnPlateau o współczynniku 0.5 i cierpliwości 5 epok na stracie testowej, seed 42); README.md repozytorium klasyfikatora, sekcje Data splits oraz Run artifacts and checkpoints (punkt kontrolny o najwyższej dokładności testowej, cel przeszukiwania hiperparametrów: dokładność testowa); pole mlp w plikach experiment.json.

\paragraph{Model szumu i jego poziomy.}
W symulacjach stosowano model warstwowy z podrozdziału~\ref{subsec:tb:noise:models}, bez błędu odczytu: po każdym bloku na każdy kubit działają odstrojenie $R_z(\cohangle)$, tłumienie fazy o parametrze $\qphase$ i kanał depolaryzujący o parametrze $\qdep$ (tabela~\ref{tab:tb:noise:channels}).
Szum działa przed końcową warstwą bramek Hadamarda, która zamienia osie $x$ i $z$, więc w modelu jednoblokowym zmierzony wektor Blocha $\blochvec_j'$ jest obrazem wektora bez szumu w odwzorowaniu postaci~\eqref{eq:tb:noise:affine}:
\begin{equation}\label{eq:cls:noise-frame}
r_{j,x}'=\distshrink r_{j,x},\quad
r_{j,y}'=\distshrink c_{\varphi}(r_{j,y}\cos\cohangle-r_{j,z}\sin\cohangle),\quad
r_{j,z}'=\distshrink c_{\varphi}(r_{j,y}\sin\cohangle+r_{j,z}\cos\cohangle),
\end{equation}
gdzie $\distshrink=1-\tfrac{4}{3}\qdep$ i $c_{\varphi}=\sqrt{1-\qphase}$.
Depolaryzacja nie zmienia zatem azymutu, tłumienie fazy przyciąga go do osi $x$, leżącej w połowie między kierunkami wzorcowymi, a odstrojenie jest obrotem wokół osi $x$ i zmienia azymut w stopniu zależnym od składowej $r_{j,z}$; przy większej liczbie bloków szum wcześniejszych bloków przekształcają kolejne bramki dwukubitowe i odwzorowanie nie jest już lokalne.

Siły składowych zadawano na dwa sposoby: jako poziomy stopniowane $0$, $0{,}01$, $0{,}05$ i $0{,}10$, używane w zestawieniu symulatorowym, oraz jako poziom ibm\_fez, wyznaczony z kalibracji stanami GHZ, $(\ket{0}^{\otimes n}+\ket{1}^{\otimes n})/\sqrt{2}$.
Stany te przygotowano na urządzeniu ibm\_fez dla $n$ od $2$ do $7$ i zmierzono z $8192$ powtórzeniami każdego obwodu, bez mitygacji błędów; wierność wyznaczono metodą koherencji wielokwantowych~\cite{wei_verifying_2020}, a trójkę $(\cohangle,\qphase,\qdep)$ dobrano tak, by symulacja obwodu przygotowującego stan, z szumem po początkowej bramce Hadamarda i po każdej bramce CNOT, odtwarzała zmierzoną wierność.
Wierność wynosi od $0{,}985$ dla dwóch kubitów do $0{,}863$ dla siedmiu, a dopasowane siły podaje tabela~\ref{tab:cls:models}: modele M1, M2 i M6 uczono z trójką dla $n=6$, pozostałe z trójką dla własnej liczby kubitów.
% ŹRÓDŁO: torus_classifier/config.py (GRADED_NOISE_LEVELS, IBM_FEZ_GHZ6_CALIBRATION) i torus_classifier/ibm_fez_ghz_calibrations.json; opis procedury: manuskrypt, sekcja 3.5 (część napisana przez autora rozprawy).
W obwodach GHZ dawka szumu następuje po każdej bramce CNOT, a w klasyfikatorze po całym bloku zawierającym $2n(n-1)$ bramek dwukubitowych, więc poziom ibm\_fez nie jest modelem urządzenia dla obwodów klasyfikatora, lecz punktem odniesienia na osi siły szumu.

\paragraph{Przebiegi na urządzeniu.}
Zapisane wyniki pochodzą z procesora ibm\_fez ($156$ kubitów, rodzina Heron~r2), z przebiegów wykonanych od września do listopada 2026~r. bez mitygacji błędów, losowej kompilacji bramek (ang. \english{twirling}) i dynamicznego odsprzęgania.
W przebiegu testowym losowano z części testowej, warstwowo względem klas, $100$ próbek; dla każdej z nich obwód o kątach wyznaczonych przez nauczoną sieć odtwarzano w Qiskit~\cite{javadi-abhari_quantum_2024} bez operacji modelu szumu, transpilowano na poziomie optymalizacji~3 na ustalone kubity fizyczne i wykonywano w trzech bazach pomiarowych, $\sigma_x$, $\sigma_y$ i $\sigma_z$, po $1024$ pomiary.
Te same próbki oceniano na symulatorze z szumem uczenia, na symulatorze bez szumu i, z wyjątkiem modelu M1, siecią MLP.
Modele M2--M11 poddano następnie dostrajaniu na urządzeniu w trzech rundach, z których każdą otwiera pomiar $400$ próbek uczących w protokole przebiegu testowego.
Pierwszy z tych pomiarów, nazywany dalej przebiegiem kalibracyjnym, wykonano tym samym, jeszcze niezmienionym modelem co przebieg testowy i na tych samych kubitach fizycznych, dlatego może on służyć do kalibracji korekty odczytu ocenianej na próbkach testowych.
% ŹRÓDŁO: torus_classifier/hardware.py i hardware_training.py (stałe DEFAULT_*); pliki experiment.json przebiegów (n_samples, shots, opt_level, backend_info, sampler_options, transpile); zgodność modelu i kubitów obu przebiegów: eksperymenty/pdu-ibm-fez/klasyfikator/wyniki/metadane.json, klucz kontrole_spojnosci.

\begin{table}[tb]
\centering
\caption[Modele klasyfikatora użyte w przebiegach na urządzeniu i w eksperymentach z korektą odczytu]{Modele klasyfikatora uruchomione na urządzeniu, w kolejności malejącej liczby bramek dwukubitowych CZ po transpilacji. $K$: liczba klas, $d$: liczba cech, $n$: liczba kubitów, $L$: liczba bloków; $\cohangle$ i $\qdep$: poziom ibm\_fez, z którym model uczono ($\qphase=0$, w modelu M10 $1{,}5\cdot10^{-4}$). ZP i S: udział w eksperymencie na zapisanych pomiarach i w eksperymencie symulacyjnym; (u) oznacza analizę uzupełniającą.}
\label{tab:cls:models}
\footnotesize
\setlength{\tabcolsep}{5pt}
\begin{tabular}{@{}llrrrrrrrcc@{}}
\toprule
Model & Identyfikator & $K$ & $d$ & $n$ & $L$ & CZ & $\cohangle$ [mrad] & $\qdep$ [$10^{-3}$] & ZP & S \\
\midrule
M1 & \texttt{c7\_f200\_q6\_b15} & 7 & 200 & 6 & 15 & 1814 & 1{,}28 & 2{,}98 & (u) & tak \\
M2 & \texttt{c7\_f200\_q3\_b13\_d2} & 7 & 200 & 3 & 13 & 233 & 1{,}28 & 2{,}98 & tak & tak \\
M3 & \texttt{c6\_f300\_q5\_b3} & 6 & 300 & 5 & 3 & 198 & 5{,}98 & 4{,}33 & tak & tak \\
M4 & \texttt{c8\_f200\_q7\_b1} & 8 & 200 & 7 & 1 & 138 & 5{,}66 & 3{,}63 & tak & -- \\
M5 & \texttt{c8\_f300\_q7\_b1} & 8 & 300 & 7 & 1 & 138 & 5{,}66 & 3{,}63 & tak & -- \\
M6 & \texttt{c7\_f200\_q6\_b1} & 7 & 200 & 6 & 1 & 113 & 1{,}28 & 2{,}98 & tak & tak \\
M7 & \texttt{c6\_f100\_q5\_b1} & 6 & 100 & 5 & 1 & 63 & 5{,}98 & 4{,}33 & tak & -- \\
M8 & \texttt{c6\_f300\_q5\_b1} & 6 & 300 & 5 & 1 & 63 & 5{,}98 & 4{,}33 & tak & -- \\
M9 & \texttt{c5\_f300\_q4\_b1} & 5 & 300 & 4 & 1 & 29 & 3{,}94 & 5{,}25 & tak & -- \\
M10 & \texttt{c4\_f400\_q3\_b1} & 4 & 400 & 3 & 1 & 15 & 3{,}70 & 3{,}81 & tak & tak \\
M11 & \texttt{c3\_f200\_q2\_b1} & 3 & 200 & 2 & 1 & 3 & 9{,}02 & 4{,}89 & tak & tak \\
\bottomrule
\end{tabular}
\end{table}
% ŹRÓDŁO tabeli: pliki experiment.json przebiegów testowych (config: n_classes, n_features, n_qubits, n_blocks, deltas_per_qubit, noise_strengths; transpile.two_qubit_gates) w katalogach wymienionych w eksperymenty/pdu-ibm-fez/klasyfikator/wyniki/metadane.json i punkty_odniesienia.json; udział w eksperymentach: pliki README.md obu eksperymentów.

\paragraph{Eksperymenty z korektą odczytu.}
Na potrzeby rozprawy wykonano dwa eksperymenty, oba bez nowych przebiegów na urządzeniu i bez ponownego uczenia modeli.
W eksperymencie na zapisanych pomiarach korektę odczytu dopasowano do par wektorów Blocha z przebiegu kalibracyjnego, zmierzonego i symulowanego z szumem uczenia, a oceniono na $100$ próbkach przebiegu testowego modeli M2--M11; model M1 nie ma przebiegu kalibracyjnego i występuje tylko w analizie uzupełniającej.
Przebieg testowy poprzedzał kalibracyjny o czas od kilkunastu minut do pięciu i pół doby, więc założenie stabilności błędu~(\ref{itm:pdu:Z4}) jest w tych danych spełnione w różnym stopniu.
W eksperymencie symulacyjnym dla sześciu modeli wyznaczono na symulatorze macierzy gęstości dokładność w funkcji siły szumu: każdej składowej osobno, na siatce dziesięciu wartości od $0{,}005$ do $0{,}5$ i na poziomie ibm\_fez, oraz wszystkich trzech łącznie, jako krotności od $1$ do $50$ poziomu ibm\_fez, dla wartości dokładnych oraz dla $8192$ i $1024$ pomiarów na bazę.
Korektę kalibrowano na $400$ próbkach uczących, a oceniano na całej części testowej, w dziesięciu powtórzeniach z nowym podzbiorem kalibracyjnym i nową realizacją szumu próbkowania; odniesieniem był symulator bez szumu, a odpowiednikiem dryfu urządzenia korekta skalibrowana przy innej sile szumu niż ta, przy której jest stosowana.
W obu eksperymentach korekta nie korzysta z etykiet klas, jej parametry swobodne wybiera walidacja krzyżowa wewnątrz próbek kalibracyjnych, a wynik sprawdzają dwie kontrole negatywne: korekta dopasowana do par o losowo przestawionym przyporządkowaniu oraz ocena na pomiarach przestawionych losowo między próbkami.
% ŹRÓDŁO: eksperymenty/pdu-ibm-fez/klasyfikator/README.md (sekcje Dane, Metoda, Czas pomiarów) i eksperymenty/krzywe-degradacji/klasyfikator/README.md (sekcje Modele i dane, Metoda); w modelu M1 eksperymentu symulacyjnego podzbiór kalibracyjny jest ustalony (400 próbek).

\paragraph{Eksperyment porównawczy.}
Trzeci eksperyment porównuje korektę odczytu z treningiem w szumie i usuwa część zagrożeń trafności wymienionych niżej.
Konfigurację przeglądu szumu ($K=3$, $n=2$, $L=7$) i model M2 uczono od nowa, z ich dotychczasowymi hiperparametrami ($35$ i $75$ epok), bez szumu i na trzech poziomach stopniowanych, z wszystkimi składowymi na tym samym poziomie.
Dla każdego z pięciu ziaren zbiór dzieli się losowo na część uczącą ($3000$ próbek), walidacyjną ($1000$) i testową ($1000$), a parametry standaryzacji liczy się tylko na części uczącej; ziarno wyznacza też wagi początkowe i kolejność próbek, więc modele jednego ziarna różnią się wyłącznie szumem w uczeniu.
Harmonogram współczynnika uczenia i punkt kontrolny wybiera się na części walidacyjnej, a część testowa służy tylko do oceny wyuczonych modeli przy każdym poziomie szumu testu, przy każdej składowej osobno i, w modelu M2, przy poziomie ibm\_fez.
Korektę odczytu w wariantach FZ i AR (podrozdział~\ref{subsec:cls:results:pdu}) kalibruje się bez etykiet na próbkach uczących, a odniesieniem jest ten sam model w szumie, w którym go uczono, na wartościach dokładnych.
Odczyt ocenia się dla wartości dokładnych i dla $1024$ pomiarów na bazę, w dwudziestu realizacjach szumu próbkowania, a w modelu M2 także dla $8192$ pomiarów.
Porównania są sparowane po ziarnach: podaje się średnią różnicę z $95$-procentowym przedziałem ufności z rozkładu t~Studenta i test t o czterech stopniach swobody, a poprawkę Holma stosuje się w rodzinie pięciu porównań głównych, ustalonej przed obliczeniami; pozostałe porównania mają charakter opisowy.
Klasycznym punktem odniesienia jest sieć MLP o tych samych warstwach ukrytych, uczona i wybierana na tych samych podziałach.
% ŹRÓDŁO: eksperymenty/trening-w-szumie/OPIS.md, IMPLEMENTACJA.md i README.md oraz gleboka/OPIS.md, IMPLEMENTACJA.md i README.md (ziarna 42--46; podział z permutacji torch.randperm z ziarnem, części 0,6 / 0,2 / 0,2; korekty a2 = FZ i b2 = AR kalibrowane na 3000 próbkach uczących, analiza czułości na 1000 walidacyjnych; 20 realizacji szumu próbkowania; uzupełniająco test permutacyjny ze zmianą znaku i test McNemara w obrębie ziarna). Kod modułów klasyfikatora zmieniono tylko w podziale danych i wyborze punktu kontrolnego (roznice_wzgledem_oryginalu.diff); niezmieniona kopia kodu odtwarza oryginalne przebiegi co do bitu (kontrole w obu katalogach).

\paragraph{Miary i analiza statystyczna.}
Podstawową miarą jest dokładność (równanie~\eqref{eq:tb:ml:metrics}), a uzupełniają ją zgodność predykcji, czyli odsetek próbek, którym dwa źródła wyników (na przykład urządzenie i symulator) przypisują tę samą klasę, odległość euklidesowa wektorów Blocha z dwóch źródeł uśredniona po próbkach i kubitach, średni promień i średni bezwzględny błąd azymutu.
W eksperymencie symulacyjnym skuteczność korekty wyraża odzysk, czyli stosunek przyrostu dokładności po korekcie do jej spadku wywołanego szumem, podawany dla spadków nie mniejszych niż $0{,}02$.
Porównania są sparowane, na tych samych próbkach, z przedziałami ufności $95\%$ wyznaczanymi metodą bootstrap~\cite{efron_bootstrap_1979}; w eksperymencie na zapisanych pomiarach stosuje się dokładny test McNemara~\cite{mcnemar_note_1947} dla dokładności i zgodności, test Wilcoxona~\cite{bib:statistics_wilc} dla odległości oraz poprawkę Holma w rodzinie dwudziestu testów głównych.
Przy $100$ próbkach błąd standardowy dokładności sięga $0{,}05$, więc w danych z urządzenia wykrywalne są tylko różnice rzędu $0{,}1$.
\agenttask{Dodać cytowanie poprawki Holma (S. Holm, A simple sequentially rejective multiple test procedure, Scandinavian Journal of Statistics 6, 1979, s. 65--70): pozycja nie ma numeru DOI, a strona JSTOR nie pozwala na automatyczną weryfikację, więc wpisu nie utworzono; do potwierdzenia przez autora albo zastąpienia podręcznikiem z bazy.}

\paragraph{Zagrożenia trafności.}
Istniejące wyniki nie powstały z myślą o rozprawie i mają ograniczenia wpływające na interpretację liczb z podrozdziału~\ref{sec:cls:results}.
\begin{enumerate}
\renewcommand{\theenumi}{ZT\arabic{enumi}}
\renewcommand{\labelenumi}{(\theenumi)}
\item\label{itm:cls:ZT1} \emph{Wybór modelu na zbiorze testowym.}
Harmonogram uczenia, punkt kontrolny i hiperparametry wybrano na zbiorze testowym, a parametry standaryzacji policzono z udziałem jego próbek; ten \glslink{data-leakage}{przeciek danych} (podrozdział~\ref{subsec:tb:ml:supervised}) zawyża dokładności bezwzględne, także sieci MLP.
\item\label{itm:cls:ZT2} \emph{Jedno ziarno i jeden zbiór.}
Każdy z modeli M1--M11 ma jeden przebieg uczenia na jednym zbiorze, a średnie po $50$ zbiorach w zestawieniu symulatorowym opisują rozrzut wzdłuż gradientu trudności, nie niepewność estymacji.
\item\label{itm:cls:ZT3} \emph{Mała próba z urządzenia.}
Wyniki z urządzenia opierają się na $100$ próbkach testowych na model, wspólnych dla modeli uczonych na tym samym zbiorze.
\item\label{itm:cls:ZT4} \emph{Model szumu.}
Dawka szumu na blok nie zależy od liczby bramek w bloku, model nie zawiera błędu odczytu, relaksacji ani przesłuchu, a kalibrację stanami GHZ wykonano w innym dniu niż przebiegi klasyfikatora i, w dziewięciu z dziesięciu modeli M2--M11, na innych kubitach fizycznych.
\item\label{itm:cls:ZT5} \emph{Kalibracja korekty.}
Korektę odczytu kalibruje się na próbkach, które sieć widziała w uczeniu, a jej odniesieniem jest symulator, nie wartości idealne.
\end{enumerate}
Porównania sparowane, z korektą i bez niej na tych samych próbkach tego samego modelu, są na zagrożenie (\ref{itm:cls:ZT1}) mało wrażliwe, natomiast dokładności bezwzględne i porównania z siecią MLP należy czytać z tym zastrzeżeniem.
Eksperyment porównawczy usuwa w obu swoich konfiguracjach przeciek danych z zagrożenia (\ref{itm:cls:ZT1}), bo punkt kontrolny i standaryzację wyznacza bez części testowej, oraz zastępuje jeden przebieg uczenia pięcioma ziarnami (\ref{itm:cls:ZT2}); budowa sieci i liczba bloków pochodzą jednak nadal ze strojenia na zbiorze testowym, a każdą konfigurację uczy się na jednym zbiorze.
\agenttaskDONE{Wykonane 2026-11-27 (akapit Eksperyment porównawczy i zdanie o zagrożeniach ZT1 i ZT2 powyżej). Treść zadania: po wykonaniu eksperymentu E2 dopisać akapit o jego metodyce w zakresie faktycznie zastosowanym i zaktualizować zagrożenia ZT1 i ZT2, wskazując, które z nich ten eksperyment usuwa i dla jakich modeli.}
\agenttask{Po wykonaniu eksperymentu E4 (klasyfikacja obrazów po redukcji błędów; nierozpoczęty) dopisać do akapitu o danych opis zbioru obrazów, a do metodyki opis potoku LPIQE, PDU i pix2pix oraz rozłączności danych użytych do uczenia sieci pix2pix i do testu klasyfikatora.}

\section{Wyniki}\label{sec:cls:results}

\subsection{Wpływ szumu i głębokości obwodu na dokładność}\label{subsec:cls:results:noise}
\agentinfo{
Napisane w zadaniu PHD-TS-79 (2026-11-27) wyłącznie z plików wynikowych repozytorium klasyfikatora; pełne zestawienie liczb ze wskazaniem pliku i pola: \agentref{eksperymenty/weryfikacja-liczb/rozdzial-8.md} oraz rozdzial-8.csv i skrypty w katalogu rozdzial-8-skrypty (tam też generator danych rysunku i wierszy tabeli).
Liczby manuskryptu zgadzają się z plikami; testy sparowane i przedziały ufności (bootstrap, 10 000 losowań) policzono w rozprawie z zapisanych predykcji na próbkę.
Zaokrąglenia: połówki w górę (na przykład 0,3775 to 0,378; manuskrypt podaje 0,377).
Nazwy wprowadzone tutaj: przegląd szumu (64 przebiegi uczenia z katalogu noisesweep), odmiana modelu M6 z siecią typu transformer.
Ustalenia interpretacyjne: odporność odczytu na skrócenie wektora Blocha jest przedstawiona jako zgodna z pracą larose_robust_2020, a nie jako wynik nowy; zmian po dostrajaniu nie nazywa się efektem, bo żaden przedział nie wyklucza zera; grupa obu modeli wieloblokowych (14 próbek zyskało, 5 straciło, p = 0,064) jest wyodrębniona po obejrzeniu wyników.
Pominięte dla budżetu stron, dostępne w plikach i w zestawieniu liczb: eksperyment z głowicą klasyczną na zamrożonej sieci (katalog encoder_head modelu M2, liczb nie weryfikowano), skrócenie i obrót odwzorowania afinicznego na kubit (device_model.json), czas pracy procesora w dostrajaniu, porównanie bez szumu i z szumem dla K = 3 i d = 200 (0,825 wobec 0,823), zestawienie symulatorowe dla K = 7 (dwa pełne przebiegi wsadowe bez odniesienia MLP), zakres dokładności przeglądu szumu w epoce o najwyższej dokładności testowej (0,761--0,818), różnica 0,142 między modelem M1 a siecią MLP modeli M2 i M6, zmiana dokładności symulowanej na pełnym zbiorze testowym po dostrajaniu (od -0,036 do +0,010), promień około 0,05 przypisywany stanowi maksymalnie mieszanemu przy 1024 pomiarach.
Symulacja modelu M2 na fałszywym urządzeniu (fake_fez, Qiskit Aer) używa migawki kalibracyjnej dołączonej do biblioteki, a nie kalibracji z dnia pomiaru; plik uv.lock repozytorium przypina qiskit-ibm-runtime 0.49.0 (migawka z 26 lutego 2025 według podrozdziału 2.2), ale pliki przebiegu z 23 września 2026 nie zapisują wersji biblioteki, więc w tekście nie podano daty migawki ani cytowania.
Wykres i tabela: dane rysunku w plikach cls_depth.dat i cls_noise_sweep.dat (katalog figures/chp_classification), generowane skryptem rysunki.py; po zmianie danych źródłowych uruchomić skrypt ponownie.
}

Podrozdział zestawia istniejące wyniki klasyfikatora odczytane z plików wynikowych przebiegów; testy sparowane i przedziały ufności policzono na potrzeby rozprawy z zapisanych predykcji.
Modele uczono ze współczynnikiem uczenia $0{,}001$, zmniejszanym o połowę po pięciu epokach bez poprawy straty testowej, na wsadach po $32$ próbki i z ziarnem $42$; punkty kontrolne modeli M1--M11 pochodzą z epok od $16$ do $41$, w których dokładność na części uczącej wynosi już od $0{,}90$ do $0{,}99$.
% ŹRÓDŁO: /Users/krzysztofwerner/git/quantum-classifier/runs/<przebieg>/hyperparameters.json (learning_rate, batch_size, seed, scheduler, epochs) i metrics.json (best.epoch, train_accuracy) modeli M1--M11; hyperparameters.json przebiegów wsadowych (35--50 epok) i przeglądu szumu (35 epok); modele M1--M11: 50 albo 75 epok; eksperymenty/weryfikacja-liczb/rozdzial-8.csv, grupa H.

\paragraph{Dokładność na symulatorze na tle sieci MLP.}
Zestawienie symulatorowe obejmuje $16$ typów problemu ($K$ od $3$ do $6$, $d$ od $100$ do $400$) w obu wariantach bloku, czyli $32$ konfiguracje uczone osobno na każdym z $50$ zbiorów serii trudności; punktem odniesienia jest w nim sieć MLP o jednej warstwie ukrytej z $64$ neuronami, a nie sieć o warstwach klasyfikatora.
Średnia dokładność klasyfikatora po $50$ zbiorach wynosi od $0{,}42$ do $0{,}87$ i w $31$ z $32$ konfiguracji jest niższa od średniej sieci MLP, przeciętnie o $0{,}034$, a najwyżej o $0{,}075$; wyjątkiem jest wariant Pauliego dla $K=5$ i $d=300$ ($0{,}631$ wobec $0{,}624$), lepszy na $37$ z $50$ zbiorów.
Na pojedynczych zbiorach klasyfikator jest lepszy w $347$, równy w $14$ i słabszy w $1239$ z $1600$ porównań.
Odchylenie standardowe po zbiorach (od $0{,}04$ do $0{,}14$, dla sieci MLP od $0{,}02$ do $0{,}05$) opisuje rozrzut wzdłuż gradientu trudności~(\ref{itm:cls:ZT2}), a nie niepewność tych średnich.
% ŹRÓDŁO: runs/batch_*/summary.json (aggregates.test_accuracy, datasets[].test_accuracy; 48 przebiegów: 32 klasyfikatora i 16 sieci MLP) oraz dataset_00/hyperparameters.json (hidden = [64] w przebiegach batch_64_mlp_*, noise_strengths, n_blocks); parowanie po nazwie zbioru w skrypcie macierz.py; rozdzial-8.csv, grupa A. Odchylenie w plikach liczone z dzielnikiem 50 (manuskrypt nazywa je odchyleniem z próby). Wariant uniwersalny ma wyższą średnią w 8 z 16 typów.

Dla modeli M1--M11 odniesieniem jest sieć MLP o tych samych warstwach ukrytych (tabela~\ref{tab:cls:hw-results}).
Na pełnym zbiorze testowym klasyfikator jest od niej słabszy we wszystkich dziesięciu modelach M2--M11, o $0{,}011$ do $0{,}081$; model M1 nie ma takiego odniesienia.
Na podzbiorach $100$ próbek nie ustępuje jej w dwóch modelach (M9: $0{,}62$ wobec $0{,}55$, M11: $0{,}82$ wobec $0{,}80$), lecz żadna z dziesięciu różnic nie jest istotna (dokładny test McNemara, $p\geq 0{,}07$).
Lepsza od sieci MLP na pełnym zbiorze testowym jest tylko odmiana modelu M6, w której perceptron $\Theta_{\mathbf{w}}$ zastąpiono siecią typu transformer~\cite{vaswani2017attention} ($0{,}556$ wobec $0{,}520$), różniąca się od odniesienia budową sieci klasycznej.
Stwierdzenie o przewadze sieci MLP z podrozdziału~\ref{sec:cls:scope} dotyczy zatem porównań z siecią o tych samych warstwach na pełnym zbiorze testowym.
% ŹRÓDŁO: runs/<przebieg>/hardware_ibm_fez_*/experiment.json (checkpoint.test_accuracy, mlp.best_test_accuracy, mlp.hidden) i results.csv (sim_noisy_correct, mlp_correct) modeli M1--M11; runs/run_1_tf256x4h4p200dr20_torus_c7_f200_universal_noise444_01-10-2026-10-56-52/metrics.json (best.test_accuracy = 0,5555, epoka 76) i hyperparameters.json; rozdzial-8.csv, grupy B, C i G. Model M1 jest o 0,142 słabszy od sieci MLP modeli M2 i M6 (ten sam zbiór, inne warstwy).

\begin{table}[tb]
\centering
\caption[Modele M1--M11 na symulatorze i na urządzeniu ibm\_fez przed dostrajaniem i po nim]{Modele M1--M11 na symulatorze i na urządzeniu ibm\_fez. CZ: liczba bramek dwukubitowych po transpilacji. Pełny zbiór: dokładność klasyfikatora (FBS, symulator) i sieci MLP o tych samych warstwach ukrytych na $2000$ próbkach testowych. S: symulator z szumem uczenia (wartości dokładne), U: urządzenie; zg.: zgodność predykcji obu źródeł; S/U: liczba próbek poprawnych tylko na symulatorze i tylko na urządzeniu; $+/-$: liczba próbek, które po dostrajaniu stały się poprawne i błędne; $p$: $p$-wartość dokładnego testu McNemara.}
\label{tab:cls:hw-results}
\footnotesize
\setlength{\tabcolsep}{3.6pt}
\begin{tabular}{@{}lrcccccccccccc@{}}
\toprule
 & & \multicolumn{2}{c}{Pełny zbiór} & \multicolumn{5}{c}{Przebieg testowy, $100$ próbek} & \multicolumn{2}{c}{Średni promień} & \multicolumn{3}{c}{Po dostrajaniu} \\
\cmidrule(lr){3-4}\cmidrule(lr){5-9}\cmidrule(lr){10-11}\cmidrule(l){12-14}
Model & CZ & FBS & MLP & S & U & zg. & S/U & $p$ & S & U & U & $+/-$ & $p$ \\
\midrule
M1 & 1814 & $0{,}378$ & -- & $0{,}34$ & $0{,}13$ & $0{,}12$ & 29/8 & $<0{,}001$ & $0{,}69$ & $0{,}09$ & -- & -- & -- \\
M2 & 233 & $0{,}439$ & $0{,}520$ & $0{,}45$ & $0{,}35$ & $0{,}71$ & 14/4 & $0{,}03$ & $0{,}84$ & $0{,}37$ & $0{,}41$ & 11/5 & $0{,}21$ \\
M3 & 198 & $0{,}455$ & $0{,}509$ & $0{,}46$ & $0{,}43$ & $0{,}91$ & 4/1 & $0{,}38$ & $0{,}90$ & $0{,}50$ & $0{,}46$ & 3/0 & $0{,}25$ \\
M4 & 138 & $0{,}425$ & $0{,}477$ & $0{,}41$ & $0{,}42$ & $0{,}97$ & 0/1 & $1$ & $0{,}92$ & $0{,}67$ & $0{,}46$ & 6/2 & $0{,}29$ \\
M5 & 138 & $0{,}327$ & $0{,}400$ & $0{,}43$ & $0{,}41$ & $0{,}96$ & 2/0 & $0{,}50$ & $0{,}88$ & $0{,}64$ & $0{,}39$ & 3/5 & $0{,}73$ \\
M6 & 113 & $0{,}455$ & $0{,}520$ & $0{,}46$ & $0{,}47$ & $0{,}89$ & 1/2 & $1$ & $0{,}91$ & $0{,}61$ & $0{,}48$ & 5/4 & $1$ \\
M7 & 63 & $0{,}606$ & $0{,}660$ & $0{,}53$ & $0{,}53$ & $0{,}97$ & 1/1 & $1$ & $0{,}92$ & $0{,}71$ & $0{,}51$ & 5/7 & $0{,}77$ \\
M8 & 63 & $0{,}503$ & $0{,}524$ & $0{,}51$ & $0{,}50$ & $0{,}99$ & 1/0 & $1$ & $0{,}95$ & $0{,}77$ & $0{,}49$ & 3/4 & $1$ \\
M9 & 29 & $0{,}568$ & $0{,}605$ & $0{,}62$ & $0{,}61$ & $0{,}98$ & 1/0 & $1$ & $0{,}95$ & $0{,}78$ & $0{,}58$ & 3/6 & $0{,}51$ \\
M10 & 15 & $0{,}624$ & $0{,}648$ & $0{,}58$ & $0{,}58$ & $1{,}00$ & 0/0 & $1$ & $0{,}88$ & $0{,}83$ & $0{,}60$ & 6/4 & $0{,}75$ \\
M11 & 3 & $0{,}807$ & $0{,}818$ & $0{,}82$ & $0{,}82$ & $1{,}00$ & 0/0 & $1$ & $0{,}99$ & $0{,}94$ & $0{,}81$ & 0/1 & $1$ \\
\bottomrule
\end{tabular}
\end{table}
% ŹRÓDŁO tabeli: przebiegi testowe runs/<przebieg>/hardware_ibm_fez_*/ modeli M1--M11: experiment.json (transpile.two_qubit_gates, checkpoint.test_accuracy, mlp.best_test_accuracy), summary.json (sources.sim_noisy.r_mean, sources.hardware.r_mean), results.csv (kolumny sim_noisy_correct, hardware_correct, sim_noisy_pred, hardware_pred); dostrajanie: runs/<przebieg>/hwtrain_ibm_fez_*/eval/final_hardware/results.csv (hardware_correct). Katalogi przebiegów: eksperymenty/weryfikacja-liczb/rozdzial-8-skrypty/common.py. Wiersze wygenerował skrypt rysunki.py; testy McNemara policzono w rozprawie.

\paragraph{Głębokość obwodu na urządzeniu.}
Tabela~\ref{tab:cls:hw-results} i rysunek~\ref{fig:cls:noise-results}a zestawiają przebiegi testowe z symulatorem na tych samych $100$ próbkach.
W ośmiu modelach jednoblokowych, od $3$ do $138$ bramek CZ, dokładność na urządzeniu różni się od symulowanej najwyżej o $0{,}02$ i w żadnym z nich różnica nie jest istotna, a predykcje obu źródeł są identyczne dla $776$ z $800$ próbek (łączna różnica dokładności $-0{,}003$, $95$-procentowy przedział ufności $[-0{,}010;\,+0{,}005]$).
W modelach wieloblokowych różnica rośnie z liczbą bramek: $-0{,}03$ $[-0{,}08;\,+0{,}01]$ w modelu M3 ($198$ bramek CZ), $-0{,}10$ $[-0{,}18;\,-0{,}02]$ w modelu M2 ($233$) i $-0{,}21$ $[-0{,}32;\,-0{,}10]$ w modelu M1 ($1814$), który klasyfikuje na urządzeniu poprawnie $13$ próbek na $100$ przy $14$ oczekiwanych dla wyboru losowego.
Średni promień na urządzeniu maleje przy tym od $0{,}94$ do $0{,}09$.
% ŹRÓDŁO: results.csv i summary.json przebiegów testowych (jak w tabeli); przedziały bootstrap (10 000 losowań, ziarno 20261127) i zestawienie łączne ośmiu modeli jednoblokowych M4--M11 (urządzenie 0,5425, symulator 0,5450, różnica -0,0025): skrypty modele.py i dostrajanie.py; rozdzial-8.csv, grupy B i C.
Liczba bloków nie była zmieniana niezależnie od typu problemu i terminu pomiaru, ale ten sam obraz dają modele uczone na tym samym zbiorze (M6, M2 i M1 oraz M8 i M3); model jednoblokowy jest w obu grupach najlepszy także na symulatorze, więc większa głębokość nie daje zysku równoważącego stratę na urządzeniu.

\begin{figure}[tb]
\centering
% !TeX root = ./../../main.tex
% Rysunek fig:cls:noise-results: (a) średni promień wektora Blocha i zgodność predykcji w funkcji liczby bramek CZ,
% (b) dokładność a średni promień w przeglądzie szumu.
% Wykresy odtworzone w pgfplots z plików danych; żadna wartość nie została odczytana z rysunków manuskryptu.
% Dane (a): figures/chp_classification/cls_depth.dat
%   (źródło: /Users/krzysztofwerner/git/quantum-classifier/runs/<przebieg>/hardware_ibm_fez_*/experiment.json, summary.json, results.csv).
% Dane (b): figures/chp_classification/cls_noise_sweep.dat
%   (źródło: /Users/krzysztofwerner/git/quantum-classifier/runs/noisesweep_7_435_331_187_286_202_torus_c3_f200_29-07-2026-16-26-05/results.csv).
% Oba pliki danych generuje eksperymenty/weryfikacja-liczb/rozdzial-8-skrypty/rysunki.py.
% Wymaga: pgfplots z biblioteką groupplots (wczytane w config/my-settings.tex).
\begin{tikzpicture}
\begin{groupplot}[
  group style={group size=2 by 1, horizontal sep=1.7cm},
  width=0.5\textwidth,
  height=4.6cm,
  ymajorgrids,
  grid style={black!15},
  tick label style={font=\footnotesize},
  label style={font=\footnotesize},
  title style={font=\footnotesize, yshift=-1mm},
  legend style={font=\scriptsize, draw=black!40, fill=white, inner sep=1.5pt, row sep=-1.5pt},
  legend cell align=left,
  yticklabel style={/pgf/number format/.cd, fixed, fixed zerofill, precision=1, use comma},
]
% ---------- (a) modele M1--M11: promień i zgodność predykcji a liczba bramek CZ ----------
\nextgroupplot[
  title={(a) modele M1--M11},
  xmode=log,
  log ticks with fixed point,
  xmin=2, xmax=3200,
  xtick={3,10,30,100,300,1000},
  xticklabels={3,10,30,100,300,1000},
  ymin=0, ymax=1.05,
  ytick={0,0.2,0.4,0.6,0.8,1.0},
  xlabel={liczba bramek CZ po transpilacji},
  ylabel={średni promień, zgodność},
  legend pos=south west,
]
\addplot[only marks, mark=o, mark size=2.2pt, black] table[x=cz, y=rsim] {figures/chp_classification/cls_depth.dat};
\addlegendentry{promień, symulator}
\addplot[only marks, mark=*, mark size=2.2pt, black] table[x=cz, y=rhw] {figures/chp_classification/cls_depth.dat};
\addlegendentry{promień, urządzenie}
\addplot[only marks, mark=x, mark size=3pt, thick, black!75] table[x=cz, y=agree] {figures/chp_classification/cls_depth.dat};
\addlegendentry{zgodność predykcji}
\node[font=\scriptsize, anchor=east, inner sep=4pt] at (axis cs:1814,0.105) {M1};
\node[font=\scriptsize, anchor=north west, inner sep=3pt] at (axis cs:233,0.3658) {M2};
\node[font=\scriptsize, anchor=north east, inner sep=3pt] at (axis cs:198,0.5032) {M3};
% ---------- (b) przegląd szumu: dokładność a promień ----------
\nextgroupplot[
  title={(b) przegląd szumu, 64 kombinacje},
  xmin=0.15, xmax=1.05,
  xtick={0.2,0.4,0.6,0.8,1.0},
  xticklabel style={/pgf/number format/.cd, fixed, fixed zerofill, precision=1, use comma},
  ymin=0.70, ymax=0.85,
  ytick={0.70,0.75,0.80,0.85},
  yticklabel style={/pgf/number format/.cd, fixed, fixed zerofill, precision=2, use comma},
  xlabel={średni promień},
  ylabel={dokładność},
  legend pos=south west,
  legend columns=2,
  legend style={/tikz/every even column/.append style={column sep=2mm}},
]
\addplot[only marks, mark=o, mark size=2pt, black, restrict expr to domain={\thisrow{dep}}{0:0}] table[x=r, y=accfinal] {figures/chp_classification/cls_noise_sweep.dat};
\addlegendentry{$\qdep=0$}
\addplot[only marks, mark=square, mark size=1.8pt, black, restrict expr to domain={\thisrow{dep}}{1:1}] table[x=r, y=accfinal] {figures/chp_classification/cls_noise_sweep.dat};
\addlegendentry{$\qdep=0{,}01$}
\addplot[only marks, mark=triangle, mark size=2.4pt, black, restrict expr to domain={\thisrow{dep}}{2:2}] table[x=r, y=accfinal] {figures/chp_classification/cls_noise_sweep.dat};
\addlegendentry{$\qdep=0{,}05$}
\addplot[only marks, mark=*, mark size=1.8pt, black, restrict expr to domain={\thisrow{dep}}{3:3}] table[x=r, y=accfinal] {figures/chp_classification/cls_noise_sweep.dat};
\addlegendentry{$\qdep=0{,}10$}
\end{groupplot}
\end{tikzpicture}

\caption[Promień wektora Blocha i zgodność predykcji a liczba bramek oraz dokładność a promień w przeglądzie szumu]{Szum a stan odczytu i decyzja klasyfikatora. (a)~Średni promień wektora Blocha na symulatorze z szumem uczenia i na urządzeniu ibm\_fez oraz zgodność predykcji obu źródeł w funkcji liczby bramek CZ po transpilacji; modele M1--M11, po $100$ próbek testowych. (b)~Dokładność na części testowej a średni promień w ostatniej epoce uczenia dla $64$ kombinacji poziomów stopniowanych w przeglądzie szumu ($K=3$, $d=200$, $n=2$, $L=7$, wartości dokładne); znacznik oznacza siłę depolaryzacji $\qdep$. Opracowanie własne.}
\label{fig:cls:noise-results}
\end{figure}

\paragraph{Odporność odczytu kierunkowego.}
W przeglądzie szumu jedną konfigurację ($K=3$, $d=200$, $n=2$, $L=7$, wariant Pauliego) uczono od nowa dla każdej z $64$ kombinacji poziomów stopniowanych trzech składowych szumu, na zbiorze spoza serii trudności, i oceniano w szumie uczenia na wartościach dokładnych.
Dokładność na części testowej w ostatniej epoce wynosi od $0{,}758$ do $0{,}815$, choć średni promień maleje od $0{,}99$ do $0{,}25$ (rysunek~\ref{fig:cls:noise-results}b), a współczynnik korelacji obu wielkości to $-0{,}05$.
Na kolejnych poziomach depolaryzacji średnia dokładność wynosi $0{,}793$, $0{,}795$, $0{,}795$ i $0{,}799$ przy średnim promieniu $0{,}84$, $0{,}77$, $0{,}52$ i $0{,}31$; z rosnącą siłą szumu średnia dokładność maleje monotonicznie dla odstrojenia, z $0{,}802$ do $0{,}786$ przy prawie niezmienionym promieniu ($0{,}62$ i $0{,}60$), oraz słabiej dla tłumienia fazy, z $0{,}798$ do $0{,}792$ przy promieniu malejącym z $0{,}69$ do $0{,}50$.
% ŹRÓDŁO (dopisek orkiestratora, 2026-11-27): średnie po 16 przebiegach na poziom, przeliczone niezależnie z tego samego pliku results.csv; tłumienie fazy: 0,7977; 0,7973; 0,7948; 0,7920; odstrojenie: 0,8021; 0,7994; 0,7939; 0,7864.
% ŹRÓDŁO: runs/noisesweep_7_435_331_187_286_202_torus_c3_f200_29-07-2026-16-26-05/results.csv (final_test_accuracy, test_r_mean; 64 wiersze, status ok; w epoce o najwyższej dokładności testowej, kolumna best_test_accuracy, zakres wynosi 0,761--0,818), summary.json i noise000/hyperparameters.json (zbiór data/synthetic_dataset_3class_5000x200.npz, 3000 i 2000 próbek, 35 epok), noise000/architecture.txt (warstwa wyjściowa 70 = 7 bloków po 10 kątów, czyli wariant Pauliego dla n = 2); średnie po poziomach i korelacja: skrypt uczenie.py; błąd standardowy pojedynczej dokładności na 2000 próbkach: 0,009; rozdzial-8.csv, grupa D.
Jest to zgodne z równaniem~\eqref{eq:cls:noise-frame}: depolaryzacja skraca wektor Blocha bez zmiany azymutu, a decyzja według wzoru~\eqref{eq:cls:readout} zależy tylko od azymutów.
Nie jest to obserwacja nowa: LaRose i~Coyle~\cite{larose_robust_2020} udowodnili, że decyzji binarnego klasyfikatora progowego nie zmienia depolaryzacja, a przy ograniczeniu prawdopodobieństw błędów także inne kanały Pauliego.
Twierdzenia te, podobnie jak przegląd szumu, dotyczą dokładnych prawdopodobieństw; przy skończonej liczbie pomiarów skrócenie wektora zwiększa rozrzut azymutu zgodnie z oszacowaniem~\eqref{eq:cls:azimuth-noise}.
Przy $1024$ pomiarach na bazę odporność potwierdzają modele jednoblokowe na urządzeniu: średni promień wynosi w nich od $0{,}61$ do $0{,}94$, a zgodność predykcji z symulatorem od $0{,}89$ do $1{,}00$; w modelu M7, zmierzonym dwukrotnie w odstępie około piętnastu godzin na tych samych kubitach fizycznych, promień wyniósł $0{,}62$ i $0{,}71$, a wszystkie $100$ predykcji było identycznych.
% ŹRÓDŁO: przebiegi modelu M7 hardware_ibm_fez_06-11-2026-09-40-05 i hardware_ibm_fez_07-11-2026-00-27-30: summary.json (sources.hardware.r_mean = 0,6177 i 0,7089), results.csv (hardware_pred), experiment.json (physical_qubits, created_at); rozdzial-8.csv, grupa G. Twierdzenia pracy larose_robust_2020: notatka eksperymenty/przeglad-literatury/sekcja-3.7-klasyfikacja.md (Tw. 1, 3 i 5; praca czytana w pełnym tekście).

\paragraph{Niedopasowanie poziomu szumu.}
Przegląd szumu dotyczy modeli ocenianych w szumie, w którym je uczono, dlatego model M2, uczony na poziomie ibm\_fez, oceniono także bez ponownego uczenia na symulatorze z trzema składowymi szumu ustawionymi jednocześnie na $0{,}01$, $0{,}05$ i $0{,}10$.
Dokładność na $100$ próbkach spada wtedy z $0{,}45$ do $0{,}43$, $0{,}40$ i $0{,}21$, a średni promień z $0{,}84$ do $0{,}65$, $0{,}20$ i $0{,}06$.
Istotny jest tylko ostatni spadek, $-0{,}24$ $[-0{,}36;\,-0{,}12]$; dla siły $0{,}05$ różnica wynosi $-0{,}05$ $[-0{,}11;\,+0{,}01]$.
Sam promień tej straty nie tłumaczy: stany produktowe o tych samych kierunkach wektorów Blocha, przygotowane jedną warstwą obrotów jednokubitowych i poddane trzynastu dawkom szumu o sile $0{,}10$, dają dokładność $0{,}42$ przy promieniu $0{,}12$ (różnica $-0{,}03$ $[-0{,}09;\,+0{,}03]$), przy czym szum działa w nich tuż przed pomiarem, w innym układzie osi niż w klasyfikatorze.
Wskazuje to, że o stracie decyduje szum wcześniejszych bloków, który kolejne bramki dwukubitowe przekształcają w zmianę kierunku wektorów Blocha (podrozdział~\ref{sec:cls:methodology}) i do którego sieć uczona przy $\cohangle=1{,}3$~mrad i $\qphase=0$ nie mogła się dostosować.
Jest to wynik jednego modelu na jednym zbiorze~(\ref{itm:cls:ZT2}, \ref{itm:cls:ZT3}); zależność od każdej składowej osobno bada eksperyment symulacyjny (podrozdział~\ref{subsec:cls:results:pdu}).
% ŹRÓDŁO: runs/run_13_512_512_256_128_256_torus_c7_f200_d2_universal_noise444_23-09-2026-09-49-24/replay_sim_23-09-2026-22-39-53, -22-39-55 i -22-39-57: experiment.json (experiment_noise_strengths, deep_noise_layers = 13), summary.json (sources.trained_noisy, sources.trained_noisy_levels, sources.replay_noisy_deep: accuracy, r_mean; prediction_agreement.trained_noisy) i results.csv (kolumny *_correct, testy sparowane w skrypcie powtorzenia.py); opis źródeł: torus_classifier/bloch_replay.py; rozdzial-8.csv, grupa E.

\paragraph{Model szumu a urządzenie.}
Poziom ibm\_fez prawie nie zmienia stanów modeli jednoblokowych: średni promień na symulatorze z szumem jest najwyżej o $0{,}007$ mniejszy niż bez szumu, a na urządzeniu jest od symulowanego mniejszy o $0{,}04$ do $0{,}30$ (w modelach wieloblokowych o $0{,}40$ do $0{,}60$).
Model szumu zaniża zatem utratę czystości stanu, zgodnie z zagrożeniem~(\ref{itm:cls:ZT4}): dawka szumu skalibrowana na jedną bramkę CNOT przypada w nim na cały blok.
Symulacja modelu M2 w Qiskit Aer z modelem szumu na poziomie bramek, zbudowanym z migawki kalibracyjnej ibm\_fez dołączonej do biblioteki, daje średni promień $0{,}38$ wobec $0{,}37$ na urządzeniu, ale dokładność $0{,}47$ i zgodność z symulatorem uczenia $0{,}92$ wobec $0{,}35$ i $0{,}71$: odtwarza utratę czystości, lecz nie zmianę decyzji.
Odchylenie urządzenia od symulatora jest przy tym w dużej części systematyczne: odwzorowanie afiniczne wektora Blocha symulowanego w zmierzony, dopasowane dla każdego kubitu w pierwszej rundzie dostrajania, ma współczynnik determinacji $R^2$ od $0{,}91$ do $0{,}996$ w modelach M3--M11 i $0{,}78$ w modelu M2, a odwzorowanie tożsamościowe od $0{,}99$ (M10, M11) przez $0{,}38$ (M4) do wartości ujemnych (M2, M3).
% ŹRÓDŁO: summary.json przebiegów testowych (sources.sim_noiseless.r_mean, sources.sim_noisy.r_mean, sources.hardware.r_mean); runs/run_13_..._d2_universal_noise444_23-09-2026-09-49-24/hardware_fake_fez_23-09-2026-22-13-31/summary.json (sources.hardware: accuracy 0,47, r_mean 0,3790; prediction_agreement.hardware.sim_noisy 0,92) i experiment.json (backend fake_fez); runs/<przebieg>/hwtrain_ibm_fez_*/experiment.json, history[0].device_model (r2_mean, r2_identity_mean; 400 próbek uczących); rozdzial-8.csv, grupy B, E i F.

\paragraph{Dostrajanie na urządzeniu.}
Każda z trzech rund dostrajania zaczyna się pomiarem $400$ próbek uczących, wylosowanych warstwowo z części uczącej; jeśli zmierzona strata wzrosła względem ostatniego zaakceptowanego pomiaru, poprzednią aktualizację sieci cofa się, a współczynnik uczenia i liczbę epok zmniejsza o połowę.
Następnie dla każdego kubitu dopasowuje się metodą najmniejszych kwadratów odwzorowanie afiniczne wektora Blocha symulowanego w zmierzony i przez pięć epok uczy sieć na tych samych próbkach (Adam, współczynnik uczenia $10^{-4}$), licząc odczyt z obrazu wektora symulowanego w tym odwzorowaniu; po trzeciej rundzie ponownie mierzy się $100$ próbek testowych.
% ŹRÓDŁO: torus_classifier/hardware_training.py (docstring modułu, funkcje acceptance, fit_device_model, train_epochs, prepare_setup; stałe DEFAULT_ROUNDS = 3, DEFAULT_EPOCHS_PER_ROUND = 5, DEFAULT_LR = 1e-4, DEFAULT_N_TRAIN = 400) i pola rounds, epochs_per_round, lr, train w plikach experiment.json dostrajania.
Na próbkach uczących procedura działa: zmierzona strata zmalała we wszystkich dziesięciu modelach, od $1{,}2$ do $2{,}8$ razy, a z $20$ pomiarów następujących po aktualizacji odrzucono trzy.
Na próbkach testowych zysku nie widać (tabela~\ref{tab:cls:hw-results}): zmiany dokładności mieszczą się między $-0{,}03$ a $+0{,}06$ i w żadnym modelu nie są istotne ($p\geq 0{,}21$), a w ośmiu modelach jednoblokowych $31$ próbek zyskało i $33$ straciły (różnica $-0{,}003$ $[-0{,}023;\,+0{,}018]$).
Dodatni znak mają zmiany w obu modelach wieloblokowych, w których urządzenie odbiegało od symulatora, $+0{,}06$ $[-0{,}02;\,+0{,}14]$ w modelu M2 i $+0{,}03$ $[0{,}00;\,+0{,}07]$ w modelu M3 (łącznie $14$ próbek zyskało, $5$ straciło, $p=0{,}064$), lecz grupę tę wyodrębniono po obejrzeniu wyników i dane nie rozstrzygają, czy dostrajanie pomaga tam, gdzie urządzenie odbiega od symulatora.
Skalę zmienności przypadkowej pokazuje powtórzenie całej sekwencji dla modelu M7: z tej samej dokładności $0{,}53$ dostrajanie dało $0{,}48$ w pierwszym i $0{,}51$ w drugim wykonaniu.
W płytkich obwodach nie ma czego odzyskiwać, bo urządzenie zgadza się z symulatorem, a dostrajanie obniża stratę na próbkach, które sieć klasyfikuje już niemal bezbłędnie (zmierzona dokładność w pierwszej rundzie od $0{,}945$ do $0{,}9925$ w modelach M3--M11).
% ŹRÓDŁO: runs/<przebieg>/hwtrain_ibm_fez_*/experiment.json (history[].hardware.loss i accuracy, history[].decision.accepted, history[-1].final_epoch.sim_test_full_accuracy) oraz eval/final_hardware/results.csv i results.csv przebiegów testowych (hardware_correct; testy sparowane w skrypcie dostrajanie.py); osiem modeli jednoblokowych łącznie: 0,5425 przed i 0,5400 po; dokładność symulowana na pełnym zbiorze testowym zmieniła się o -0,036 do +0,010; model M7: hwtrain_ibm_fez_06-11-2026-17-45-59 i hwtrain_ibm_fez_07-11-2026-00-58-36 (hardware_evaluations[].hardware_accuracy); rozdzial-8.csv, grupy F i G.

Poza zagrożeniami (\ref{itm:cls:ZT1})--(\ref{itm:cls:ZT5}) wyniki tego podrozdziału mają dwa ograniczenia własne: przegląd szumu i ocena przy innym poziomie szumu nie obejmują szumu próbkowania, a testy podano bez poprawki na wielokrotne porównania, więc mają charakter opisowy.
\agenttaskDONE{Wykonane przez orkiestratora 2026-11-27 (zdanie w podrozdziale 8.3 doprecyzowano po potwierdzeniu pola hidden w plikach przebiegów batch_64_mlp). Treść zgłoszenia: zdanie w akapicie Dane i uczenie podrozdziału 8.3, że klasycznym punktem odniesienia jest sieć MLP o tych samych warstwach ukrytych, jest prawdziwe tylko dla modeli M2--M11; w zestawieniu symulatorowym odniesieniem jest sieć o jednej warstwie ukrytej z 64 neuronami (pliki hyperparameters.json przebiegów batch_64_mlp, pole hidden), a model M1 nie ma własnego odniesienia. Podrozdział 8.4.1 podaje to wprost; warto uzgodnić brzmienie w 8.3.}

\subsection{Korekta odczytu typu PDU}\label{subsec:cls:results:pdu}
\agentinfo{
Napisane w zadaniu PHD-TS-88 (2026-11-27) z plików wynikowych eksperymentów E0 \agentref{eksperymenty/pdu-ibm-fez/klasyfikator/README.md} i E1 \agentref{eksperymenty/krzywe-degradacji/klasyfikator/README.md}; liczby policzone dodatkowo odtwarza skrypt \agentref{eksperymenty/weryfikacja-liczb/rozdzial-8-skrypty/korekta_odczytu.py}.
Nazwy wariantów w tekście i ich odpowiedniki w plikach: FO = (a) a_pdu, FZ = (a2) a2_pdu_slepa, AO = (b) b_afin, AR = (b2) b2_afin_wprost; warianty z funkcją błędu kąta biegunowego (a_pdu_polar, a2_pdu_slepa_polar, tylko model M2) pominięto.
Wariant FO podawać wyłącznie z wynikiem kontroli, w której pomiary przestawiono między próbkami; wnioski opierać na wariantach FZ, AO i AR.
W eksperymencie E0 rodzinę 20 testów głównych tworzą warianty FO i AO (tak eksperyment zaplanowano), a FZ i AR są tam eksploracyjne; tekst mówi to wprost i nie należy tego wygładzać.
Notacja: wektor zmierzony ma znak prim jak w eq:cls:noise-frame, wynik skorygowany daszek jak w rozdziale 4, funkcja błędu fazy to makro phaseerr (litera nu, nie epsilon), liczba jej węzłów to makro nnodes (nie K), macierze wariantu AO to makra ptmblock i ptmshift z indeksem kubitu.
Policzone dodatkowo z pliku wyniki/predykcje.csv eksperymentu E0 (pliki zbiorcze tego nie zawierają): wynik zbiorczy wariantów FZ i AR na 1000 predykcji (bootstrap warstwowy po modelach, 10 000 losowań, ziarno 20261127, dokładny test McNemara) oraz liczby próbek o zmienionej poprawności w ośmiu modelach jednoblokowych.
Zaokrąglenia: połówki w górę, jak w 8.4.1.
Rozbieżność README--pliki (drobna): README eksperymentu E1 podaje, że przy depolaryzacji 0,5 i 1024 pomiarach dokładność modeli jednoblokowych zmienia się najwyżej o 0,002, a w pliku krzywe.csv spadek dla modelu M11 wynosi 0,00265; w tekście tej liczby nie użyto.
Objętość w kompilacji bez notatek (2026-11-27): około 3,7 strony (tekst 3,1; tabela 0,27; rysunek 0,33) przy planowanych około 3; kandydaci do dalszego cięcia: tabela tab:cls:pdu-recovery do dodatku (0,27), analiza uzupełniająca modelu M1 (0,13), zdanie o kalibracji na odłożonej połowie części testowej (0,08), zdanie o literaturze (0,07).
Pominięte dla budżetu stron, z liczbami zapisanymi w komentarzach źródłowych przy odpowiednich zdaniach: walidacja krzyżowa w obrębie przebiegu kalibracyjnego modelu M2 (FZ i AR podnoszą tam dokładność z 0,795 do 0,890 i 0,963 na próbkach uczących; jedyny sygnał dodatni z urządzenia, eksploracyjny), wiersze przebiegu mieszanego w tabeli, niestabilność wariantu AO między powtórzeniami kalibracji, niedopasowanie kalibracji w modelu M1, porównanie promienia z urządzeniem (krotność 11 poziomu ibm_fez w modelu M2), liczba 99,7 procent faz przy kierunkach wzorcowych, zliczenia próbek w modelach jednoblokowych.
Pominięte i niezapisane w komentarzach, dostępne w plikach: krzywe modeli M1 i M3, wyniki dla 8192 pomiarów i wartości dokładnych poza zdaniami w tekście, pełne macierze niedopasowania (także dla tłumienia fazy i depolaryzacji), czułość na liczbę węzłów, druga para przebiegów modelu M7, skutki uboczne korekty przy braku spadku dokładności, kontrola z przestawionymi parami w eksperymencie E1.
Kontekst literaturowy: wyłącznie prace oznaczone w notatce \agentref{eksperymenty/przeglad-literatury/sekcja-3.7-klasyfikacja.md} jako czytane w pełnym tekście (peters_machine_2021, singh_benchmarking_2026), bez liczb; zdanie o normalizacji wyników pomiaru (wang_quantumnat_2022) usunięto dla budżetu stron.
Etykiety dodane: eq:cls:corrections, tab:cls:pdu-recovery, fig:cls:pdu-curves.
Rysunek: dane w plikach cls_pdu_m2_*.csv (katalog figures/chp_classification) są niezmienionymi kopiami plików wyniki/pgfplots eksperymentu E1 dla modelu M2; po ponownym uruchomieniu eksperymentu skopiować je ponownie.
}

\paragraph{Warianty korekty odczytu.}
Niech $\blochvec_j'$ i $\varphi_j'$ oznaczają zmierzony wektor Blocha $j$-tego kubitu i jego azymut, a $\blochvec_j$ i $\varphi_j$ wektor i azymut odniesienia z klasycznej symulacji tego samego obwodu (z szumem uczenia w eksperymencie na zapisanych pomiarach, bez szumu w symulacyjnym).
Korektę dopasowuje się dla każdego kubitu do par $(\blochvec_j',\blochvec_j)$ próbek kalibracyjnych w czterech wariantach:
\begin{equation}\label{eq:cls:corrections}
\begin{aligned}
\text{FO:}\quad&\hat{\varphi}_j=\varphi_j'-\phaseerr_j^{\mathrm{o}}(\varphi_j), &\qquad \text{AO:}\quad&\hat{\blochvec}_j=\ptmblock_j^{-1}\bigl(\blochvec_j'-\ptmshift_j\bigr),\\
\text{FZ:}\quad&\hat{\varphi}_j=\varphi_j'-\phaseerr_j^{\mathrm{z}}(\varphi_j'), &\qquad \text{AR:}\quad&\hat{\blochvec}_j=\mathbf{B}_j\blochvec_j'+\mathbf{b}_j.
\end{aligned}
\end{equation}
W wariantach fazowych funkcja błędu fazy $\phaseerr_j$ jest odpowiednikiem funkcji PDU kubitu (definicja~\ref{def:pdu:function}): jej wartością w każdym z $\nnodes$ równoodległych węzłów jest ważona średnia kołowa błędów $\varphi_j'-\varphi_j$ próbek kalibracyjnych o argumencie najbliższym temu węzłowi, między węzłami interpoluje się liniowo, a liczbę $\nnodes\leq 32$ wybiera walidacja krzyżowa; argumentem jest faza oczekiwana $\varphi_j$ (FO) albo zmierzona $\varphi_j'$ (FZ).
W wariancie AO dopasowuje się metodą najmniejszych kwadratów odwzorowanie~\eqref{eq:tb:noise:affine}, $\blochvec_j'\approx\ptmblock_j\blochvec_j+\ptmshift_j$, i odwraca je, a w wariancie AR dopasowuje się od razu regresję wektora odniesienia na wektor zmierzony.
% ŹRÓDŁO: eksperymenty/pdu-ibm-fez/klasyfikator/README.md (sekcja Metoda) i eksperymenty/krzywe-degradacji/klasyfikator/korekta.py (funkcje pdu_fit, pdu_eval, select_K, affine, fit, apply; K_GRID, N_FOLDS = 5). Szczegóły pominięte w tekście: liczbę węzłów wybiera się ze zbioru 1, 2, 4, 8, 16, 32; błąd azymutu sprowadza się do przedziału (-pi, pi]; węzły leżą w punktach -pi/2 + 2 pi k / N_w; wagą jest iloczyn długości rzutów obu wektorów na płaszczyznę xy; uśrednia się i interpoluje wskazy, nie kąty; do każdego węzła dodaje się jedną pseudoobserwację o globalnym błędzie kubitu; stratą walidacji krzyżowej jest ważone 1 - cos reszty fazy; w wariancie AO używa się pseudoodwrotności; warianty fazowe obracają wektor wokół osi z, więc nie zmieniają jego długości ani kąta biegunowego.

W chwili korekty warianty FZ, AO i AR wymagają tylko pomiaru, a wariant FO fazy oczekiwanej ocenianej próbki, czyli symulacji każdego ocenianego obwodu; ponieważ decyzja~\eqref{eq:cls:readout} jest funkcją samych azymutów, argument ten zawiera predykcję odniesienia i wprowadza do korekty informację o poprawnej odpowiedzi.
Wariant FO ma postać wyniku skorygowanego z definicji~\ref{def:pdu:correction}, lecz nie spełnia założenia~(\ref{itm:pdu:Z7}), a wariant FZ zastępuje znany argument zmierzonym, co jest przypadkiem wyłączonym w podrozdziale~\ref{par:pdu:known-argument}; w obu odniesieniem jest symulator (wbrew założeniu~\ref{itm:pdu:Z3}), a węzły nie tworzą zaprojektowanej siatki, bo fazy próbek kalibracyjnych skupiają się przy kierunkach wzorcowych.
% ŹRÓDŁO: eksperymenty/pdu-ibm-fez/klasyfikator/wyniki/metadane.json, pole modele[].udzial_faz_sym_dalej_niz_30st_od_pm90.uczace (najwyżej 0,0029 w dziesięciu modelach, czyli ponad 99,7% faz oczekiwanych próbek kalibracyjnych leży nie dalej niż 30 stopni od kierunków wzorcowych).
Wyniki wariantu FO podaje się dalej wyłącznie z wynikiem kontroli, w której pomiary przestawiono między próbkami, a wnioski opierają się na pozostałych wariantach.

\paragraph{Zapisane pomiary z urządzenia.}
Rodzinę dwudziestu testów głównych tworzą warianty FO i AO w modelach M2--M11; warianty FZ i AR zdefiniowano przed obejrzeniem wyników na próbkach testowych, lecz poza tą rodziną, więc ich wyniki są w tym eksperymencie eksploracyjne.
Żaden wariant nie zmienia dokładności, która na $1000$ predykcji wynosi przed korektą $0{,}512$: zmiana z $95$-procentowym przedziałem ufności to $0{,}000$ $[-0{,}005;\,+0{,}005]$ (FO), $+0{,}001$ $[-0{,}013;\,+0{,}015]$ (AO), $+0{,}003$ $[-0{,}003;\,+0{,}009]$ (FZ) i $+0{,}006$ $[-0{,}003;\,+0{,}016]$ (AR), a $p$-wartości testów głównych dla dokładności są po poprawce Holma równe jedności.
% ŹRÓDŁO: eksperymenty/pdu-ibm-fez/klasyfikator/wyniki/zbiorczo.csv (warianty a_pdu i b_afin: dokl_przed, dokl_delta, dokl_delta_ci_lo, dokl_delta_ci_hi) i per_model.csv (dokl_p_mcnemar_holm20 = 1 w 20 wierszach); warianty FZ i AR policzono dodatkowo z pliku predykcje.csv (kolumny etykieta, pred_sprzet, pred_a2_pdu_slepa, pred_b2_afin_wprost) skryptem eksperymenty/weryfikacja-liczb/rozdzial-8-skrypty/korekta_odczytu.py: FZ 0,515 (7 zysków, 4 straty, p = 0,55), AR 0,518 (15 zysków, 9 strat, p = 0,31); bootstrap warstwowy po modelach, 10 000 losowań, ziarno 20261127; ten sam skrypt odtwarza wiersze pliku zbiorczo.csv dla FO i AO z dokładnością do losowości bootstrapu (krańce przedziałów różnią się najwyżej o 0,001).
W ośmiu modelach jednoblokowych, w których urządzenie zgadza się z symulatorem w $776$ z $800$ predykcji (podrozdział~\ref{subsec:cls:results:noise}), nie ma czego korygować: warianty FO, FZ i AR zmieniają tam poprawność najwyżej trzech predykcji.
% ŹRÓDŁO: policzone dodatkowo z pliku predykcje.csv (modele M4--M11) skryptem korekta_odczytu.py: zyski i straty FO 0/0, FZ 0/1, AR 1/2, AO 8/15.
W modelach M6 i M2 wariant AO jest źle uwarunkowany (wskaźnik uwarunkowania macierzy $\ptmblock_j$ sięga $31$ i $93$) i obniża zgodność predykcji z symulatorem z $0{,}89$ do $0{,}64$ oraz z $0{,}71$ do $0{,}52$ ($p=0{,}0002$ i $p=0{,}049$ po poprawce Holma).
% ŹRÓDŁO: per_model.csv, wiersze b_afin (zgod_przed, zgod_po, zgod_p_mcnemar_holm20; w ośmiu pozostałych modelach wariant zmniejsza odległość wektorów Blocha od symulatora o 13--52%, kolumny odl_przed i odl_po, bez zmiany dokładności); metadane.json, pola modele[].b_afin_cond_A (największe: 31,06 i 93,21) oraz sym_min_odch_std_PCA_na_kubit.uczace (0,0137--0,0392, czyli tyle co niepewność statystyczna składowej, sigma_strzalowa_srednia 0,025--0,030: wektory odniesienia próbek kalibracyjnych leżą prawie na jednej osi i jedna kolumna macierzy jest wyznaczona przez szum próbkowania).

W modelu M2, jedynym z przebiegiem kalibracyjnym, w którym urządzenie odbiega od symulatora ($0{,}35$ wobec $0{,}45$), warianty FZ i AR podnoszą dokładność o $0{,}04$ i $0{,}05$, z przedziałami ufności $[-0{,}02;\,+0{,}10]$ i $[-0{,}03;\,+0{,}13]$, czego przy stu próbkach nie da się odróżnić od zera ($p=0{,}34$ i $p=0{,}33$).
Przebieg kalibracyjny tego modelu wykonano pięć i pół doby po testowym, a średni błąd fazy pierwszego kubitu wynosi w nich $-9{,}3^{\circ}$ i $-26{,}0^{\circ}$, więc założenie stabilności~(\ref{itm:pdu:Z4}) nie jest spełnione.
% ŹRÓDŁO: eksperymenty/pdu-ibm-fez/klasyfikator/wyniki/uzupelniajace.csv, model c7_f200_q3_b13_d2, analiza test (warianty a2_pdu_slepa i b2_afin_wprost: dokl_delta, przedział, dokl_p_mcnemar); pominięte w tekście: analiza cv_w_przebiegu_uczacym tego modelu (zagnieżdżona walidacja krzyżowa na 400 próbkach uczących tej samej sesji: dokładność 0,795 rośnie do 0,89 dla FZ i 0,9625 dla AR, przyrosty +0,095 [+0,065; +0,125] i +0,1675 [+0,125; +0,210], a wariant AO obniża ją do 0,6175; są to próbki znane sieci z uczenia); metadane.json, pola modele[0].odstep_test_kalibracja_min (7851,7--7860,9 min) i blad_globalny_fazy_st (uczace -9,31; testowe -26,0 dla kubitu 0).

Korekta dopasowana do par o przestawionym przyporządkowaniu obniża dokładność w każdym modelu i wariancie (do $0{,}12$--$0{,}35$), a po przestawieniu pomiarów między próbkami zgodność z symulatorem nie rośnie po korekcie o więcej niż $0{,}01$, z wyjątkiem modelu M2 (z $0{,}19$ do $0{,}23$--$0{,}25$ w wariantach FO, FZ i AR).
W modelu M1, w którym z braku przebiegu kalibracyjnego korektę dopasowano w walidacji krzyżowej na stu próbkach testowych, wariant FO podnosi dokładność z $0{,}13$ do $0{,}26$, a zgodność z symulatorem z $0{,}12$ do $0{,}66$, lecz zgodność $0{,}66$ daje także dla pomiarów przestawionych: poprawa pochodzi z fazy oczekiwanej podanej jako argument, a warianty FZ i AR niczego nie odzyskują ($0{,}17$ i $0{,}15$; $p\geq 0{,}58$).
% ŹRÓDŁO: eksperymenty/pdu-ibm-fez/klasyfikator/wyniki/kontrole.csv (kontrola tasowane_pary: dokl_srednia 0,12--0,349 wobec 0,35--0,82 bez korekty; kontrola tasowany_sprzet: zgod_srednia wariantu bez_korekty i wariantów korekty, różnice od -0,03 do +0,005 w dziewięciu modelach, w modelu c7_f200_q3_b13_d2 0,193 wobec 0,232; 0,246 i 0,234; model c7_f200_q6_b15: 0,655 dla a_pdu, 0,340 dla a2_pdu_slepa, 0,360 dla b2_afin_wprost przy zgodności dla pomiarów prawdziwych 0,66; 0,33 i 0,36) i uzupelniajace.csv, analiza cv_w_100_testowych_modelu_15_blokow (a_pdu: +0,13 [+0,03; +0,23]; średni promień na urządzeniu 0,09 według punkty_odniesienia.json).

\begin{table}[tb]
\centering
\caption[Korekta odczytu na progu spadku dokładności w eksperymencie symulacyjnym]{Korekta odczytu na progu spadku dokładności w eksperymencie symulacyjnym (modele wieloblokowe, $\nshots=1024$). Próg: siła składowej, przy której dokładność bez korekty spada o $0{,}05$; poziom: pierwszy poziom siatki, na którym spadek osiąga $0{,}05$, i średni promień wektora Blocha na tym poziomie. Dokładność bez szumu obwodu wynosi $0{,}377$ (M1), $0{,}437$ (M2) i $0{,}453$ (M3); w nawiasach odzysk w procentach.}
\label{tab:cls:pdu-recovery}
\footnotesize
\setlength{\tabcolsep}{4pt}
\begin{tabular}{@{}llrrrcccc@{}}
\toprule
 & & & & & & \multicolumn{3}{c}{Dokładność po korekcie (odzysk)} \\
\cmidrule(l){7-9}
Model & Składowa & Próg & Poziom & Promień & Bez korekty & FZ & AO & AR \\
\midrule
M1 & odstrojenie & $0{,}107$ & $0{,}2$ & $0{,}30$ & $0{,}142$ & $0{,}218$ ($32$) & $0{,}211$ ($29$) & $0{,}319$ ($75$) \\
 & tłumienie fazy & $0{,}146$ & $0{,}2$ & $0{,}25$ & $0{,}285$ & $0{,}352$ ($73$) & $0{,}312$ ($29$) & $0{,}363$ ($85$) \\
 & depolaryzacja & $0{,}040$ & $0{,}05$ & $0{,}09$ & $0{,}290$ & $0{,}318$ ($32$) & $0{,}241$ ($-57$) & $0{,}304$ ($16$) \\
\addlinespace
M2 & odstrojenie & $0{,}067$ & $0{,}07$ & $0{,}85$ & $0{,}378$ & $0{,}384$ ($11$) & $0{,}421$ ($74$) & $0{,}427$ ($83$) \\
 & tłumienie fazy & $0{,}186$ & $0{,}2$ & $0{,}41$ & $0{,}383$ & $0{,}383$ ($-1$) & $0{,}416$ ($61$) & $0{,}412$ ($53$) \\
 & depolaryzacja & $0{,}071$ & $0{,}1$ & $0{,}07$ & $0{,}279$ & $0{,}301$ ($14$) & $0{,}220$ ($-37$) & $0{,}262$ ($-11$) \\
\addlinespace
M3 & odstrojenie & $0{,}237$ & $0{,}3$ & $0{,}90$ & $0{,}329$ & $0{,}366$ ($30$) & $0{,}454$ ($100$) & $0{,}453$ ($100$) \\
 & tłumienie fazy & \multicolumn{7}{l}{próg nieosiągnięty: do siły $0{,}5$ zmiana dokładności poniżej $0{,}002$} \\
 & depolaryzacja & $0{,}309$ & $0{,}5$ & $0{,}01$ & $0{,}192$ & $0{,}155$ ($-14$) & $0{,}181$ ($-4$) & $0{,}154$ ($-14$) \\
\bottomrule
\end{tabular}
\end{table}
% ŹRÓDŁO tabeli: eksperymenty/krzywe-degradacji/klasyfikator/wyniki/progi.csv (2000 próbek testowych, średnie z dziesięciu powtórzeń; siła progowa z interpolacji liniowej w logarytmie siły), wiersze o strzaly = 1024 modeli c7_f200_q6_b15 (M1), c7_f200_q3_b13_d2 (M2) i c6_f300_q5_b3 (M3): sila_progowa, poziom_siatki, promien_analityczny, dokl_zero_szumu, dokl_bez_korekty, dokl_a2 i odzysk_a2, dokl_b i odzysk_b, dokl_b2 i odzysk_b2; te same wartości z większą liczbą cyfr w pliku krzywe.csv (dokl, odzysk_dokl); wiersz M3, tłumienie fazy: progi.csv, spadek_przy_sile_maks = -0,0016. Zaokrąglenia połówek w górę (na przykład 0,31885 to 0,319). Pominięte dla budżetu stron wiersze przebiegu mieszanego (siła = krotność poziomu ibm_fez; próg, poziom, promień, bez korekty, FZ, AO, AR z odzyskiem w procentach): M1 12,7; 20; 0,06; 0,234; 0,260 (18); 0,201 (-23); 0,238 (3); M2 22,6; 50; 0,02; 0,169; 0,160 (-3); 0,165 (-1); 0,155 (-5); M3 28,6; 50; 0,19; 0,309; 0,378 (47); 0,446 (95); 0,452 (99).

\begin{figure}[tb]
\centering
% !TeX root = ./../../main.tex
% Rysunek fig:cls:pdu-curves: dokładność modelu M2 (c7_f200_q3_b13_d2) w funkcji siły jednej składowej szumu,
% bez korekty odczytu i z korektą w wariantach FZ i AR (eksperyment symulacyjny, 2000 próbek testowych).
% Wykres odtworzony w pgfplots z plików danych; żadna wartość nie została odczytana z rysunków.
% Dane: figures/chp_classification/cls_pdu_m2_<składowa>_<pomiary>.csv, czyli niezmienione kopie plików
%   eksperymenty/krzywe-degradacji/klasyfikator/wyniki/pgfplots/c7_f200_q3_b13_d2__<składowa>__<pomiary>.csv
%   (składowa: odstrojenie, tlumienie_fazy, depolaryzacja; pomiary: 1024 albo dokladne).
% Kolumny: sila, dokl_bez (bez korekty), dokl_a2 (wariant FZ), dokl_b2 (wariant AR); pierwszy wiersz (sila = 0)
%   to brak szumu obwodu i jest pomijany, bo oś jest logarytmiczna.
% Poziom ibm_fez modelu M2: odstrojenie 0,0012756 rad, depolaryzacja 0,0029837, tłumienie fazy 0 (kolumna ibm_fez = 1).
% Najlepsza stała predykcja: 0,1555 (wyniki/kontrole.csv, kolumna dokl_stala_predykcja).
% Wymaga: pgfplots z biblioteką groupplots i biblioteki calc (wczytane w config/my-settings.tex).
\begin{tikzpicture}
\pgfplotsset{
  clsbez/.style={black, thick, mark=*, mark size=1.5pt},
  clsexact/.style={black!45, line width=1.3pt, no marks},
  clsfz/.style={black, dashed, mark=square, mark size=1.6pt, mark options={solid}},
  clsar/.style={black, densely dashdotted, mark=triangle, mark size=2.2pt, mark options={solid}},
}
\begin{groupplot}[
  group style={group size=3 by 1, horizontal sep=0.45cm, y descriptions at=edge left},
  width=0.385\textwidth,
  height=4.5cm,
  xmode=log,
  xmin=0.001, xmax=0.7,
  xtick={0.001,0.01,0.1},
  ymin=0.05, ymax=0.47,
  ytick={0.1,0.2,0.3,0.4},
  ymajorgrids,
  grid style={black!15},
  tick label style={font=\footnotesize},
  label style={font=\footnotesize},
  title style={font=\footnotesize, yshift=-1mm},
  yticklabel style={/pgf/number format/.cd, fixed, fixed zerofill, precision=1, use comma},
  legend style={font=\scriptsize, draw=black!40, fill=white, inner sep=2pt, /tikz/every even column/.append style={column sep=3mm}},
  legend cell align=left,
  legend columns=4,
]
% ---------- (a) odstrojenie ----------
\nextgroupplot[
  title={(a) odstrojenie},
  xlabel={$\cohangle$ [rad na blok]},
  ylabel={dokładność},
  legend to name=clspdulegend,
]
\draw[black!40, thin, densely dotted] (axis cs:0.001,0.1555) -- (axis cs:0.7,0.1555);
\draw[black!55, thin, densely dashed] (axis cs:0.0012756,0.05) -- (axis cs:0.0012756,0.47);
\node[font=\scriptsize, rotate=90, anchor=south west, inner sep=1.5pt] at (axis cs:0.0021,0.06) {ibm\_fez};
\addplot[clsexact, skip coords between index={0}{1}] table[col sep=comma, x=sila, y=dokl_bez] {figures/chp_classification/cls_pdu_m2_odstrojenie_dokladne.csv};
\addlegendentry{bez korekty, wartości dokładne}
\addplot[clsbez, skip coords between index={0}{1}] table[col sep=comma, x=sila, y=dokl_bez] {figures/chp_classification/cls_pdu_m2_odstrojenie_1024.csv};
\addlegendentry{bez korekty}
\addplot[clsfz, skip coords between index={0}{1}] table[col sep=comma, x=sila, y=dokl_a2] {figures/chp_classification/cls_pdu_m2_odstrojenie_1024.csv};
\addlegendentry{wariant FZ}
\addplot[clsar, skip coords between index={0}{1}] table[col sep=comma, x=sila, y=dokl_b2] {figures/chp_classification/cls_pdu_m2_odstrojenie_1024.csv};
\addlegendentry{wariant AR}
% ---------- (b) tłumienie fazy ----------
\nextgroupplot[
  title={(b) tłumienie fazy},
  xlabel={$\qphase$},
]
\draw[black!40, thin, densely dotted] (axis cs:0.001,0.1555) -- (axis cs:0.7,0.1555);
\addplot[clsexact, skip coords between index={0}{1}] table[col sep=comma, x=sila, y=dokl_bez] {figures/chp_classification/cls_pdu_m2_tlumienie_fazy_dokladne.csv};
\addplot[clsbez, skip coords between index={0}{1}] table[col sep=comma, x=sila, y=dokl_bez] {figures/chp_classification/cls_pdu_m2_tlumienie_fazy_1024.csv};
\addplot[clsfz, skip coords between index={0}{1}] table[col sep=comma, x=sila, y=dokl_a2] {figures/chp_classification/cls_pdu_m2_tlumienie_fazy_1024.csv};
\addplot[clsar, skip coords between index={0}{1}] table[col sep=comma, x=sila, y=dokl_b2] {figures/chp_classification/cls_pdu_m2_tlumienie_fazy_1024.csv};
% ---------- (c) depolaryzacja ----------
\nextgroupplot[
  title={(c) depolaryzacja},
  xlabel={$\qdep$},
]
\draw[black!40, thin, densely dotted] (axis cs:0.001,0.1555) -- (axis cs:0.7,0.1555);
\draw[black!55, thin, densely dashed] (axis cs:0.0029837,0.05) -- (axis cs:0.0029837,0.47);
\node[font=\scriptsize, rotate=90, anchor=south west, inner sep=1.5pt] at (axis cs:0.0048,0.06) {ibm\_fez};
\addplot[clsexact, skip coords between index={0}{1}] table[col sep=comma, x=sila, y=dokl_bez] {figures/chp_classification/cls_pdu_m2_depolaryzacja_dokladne.csv};
\addplot[clsbez, skip coords between index={0}{1}] table[col sep=comma, x=sila, y=dokl_bez] {figures/chp_classification/cls_pdu_m2_depolaryzacja_1024.csv};
\addplot[clsfz, skip coords between index={0}{1}] table[col sep=comma, x=sila, y=dokl_a2] {figures/chp_classification/cls_pdu_m2_depolaryzacja_1024.csv};
\addplot[clsar, skip coords between index={0}{1}] table[col sep=comma, x=sila, y=dokl_b2] {figures/chp_classification/cls_pdu_m2_depolaryzacja_1024.csv};
\end{groupplot}
\node[anchor=north] at ($(group c2r1.south)+(0,-1.05cm)$) {\pgfplotslegendfromname{clspdulegend}};
\end{tikzpicture}

\caption[Dokładność modelu M2 w funkcji siły szumu bez korekty odczytu i z korektą]{Dokładność modelu M2 w eksperymencie symulacyjnym w funkcji siły jednej składowej szumu ($\nshots=1024$), bez korekty odczytu i z korektą w wariantach FZ i AR. Linia szara: bez korekty, wartości dokładne; linia kropkowana: najlepsza stała predykcja; pionowa linia przerywana: poziom ibm\_fez. Opracowanie własne.}
\label{fig:cls:pdu-curves}
\end{figure}

\paragraph{Symulacja: krzywe degradacji.}
Na poziomie ibm\_fez żaden z sześciu modeli nie traci dokładności (zmiana najwyżej o $0{,}001$ przy $1024$ pomiarach), a w modelach jednoblokowych M6, M10 i M11 żadna składowa o sile do $0{,}5$ nie obniża jej o więcej niż $0{,}017$: zgodnie z równaniem~\eqref{eq:cls:noise-frame} szum działa tam raz, tuż przed pomiarem, a depolaryzacja nie zmienia azymutu.
% ŹRÓDŁO: eksperymenty/krzywe-degradacji/klasyfikator/wyniki/krzywe.csv, wariant bez_korekty: wiersze ibm_fez = 1 i strzaly = 1024 (spadek_dokl od -0,0012 do +0,0012) oraz największy spadek_dokl modeli jednoblokowych (0,01675: c7_f200_q6_b1, odstrojenie 0,5, strzaly 8192; w modelach c4_f400_q3_b1 i c3_f200_q2_b1 najwyżej 0,0027).
W modelach wieloblokowych M1, M2 i M3 dokładność spada o $0{,}05$ przy siłach podanych w tabeli~\ref{tab:cls:pdu-recovery}, a krzywe modelu M2 przedstawia rysunek~\ref{fig:cls:pdu-curves}.
Liczba pomiarów ma znaczenie dopiero przy małym promieniu, zgodnie z oszacowaniem~\eqref{eq:cls:azimuth-noise}: w modelu M2 przy $\qdep=0{,}1$ (promień $0{,}07$) dokładność wynosi $0{,}394$ dla wartości dokładnych, $0{,}380$ przy $8192$ i $0{,}279$ przy $1024$ pomiarach, a bez szumu obwodu skończona liczba pomiarów zmienia ją najwyżej o $0{,}002$.
% ŹRÓDŁO: krzywe.csv, wariant bez_korekty: model c7_f200_q3_b13_d2, depolaryzacja 0,1 (dokl 0,3935; 0,37965; 0,279 dla strzaly 0, 8192 i 1024; promien_analityczny 0,073); wiersze skladowa = zero wszystkich modeli (dokl dla strzaly 0 i 1024, różnica najwyżej 0,0019).

\paragraph{Symulacja: odzysk dokładności.}
Przy odstrojeniu wariant AR odzyskuje na poziomie progowym $83\%$ (M2), $75\%$ (M1) i $100\%$ (M3) utraconej dokładności, a przedziały ufności przyrostów, $[+0{,}031;\,+0{,}066]$, $[+0{,}153;\,+0{,}202]$ i $[+0{,}100;\,+0{,}146]$, nie zawierają zera; wariant FZ odzyskuje $11\%$, $32\%$ i $30\%$ (w modelu M2 nie koryguje on kąta biegunowego, z którego korzysta odczyt tego modelu).
Zgodnie z równaniem~\eqref{eq:cls:noise-frame} odstrojenie jest obrotem wokół osi $x$, więc przesunięcie azymutu zależy od składowej $r_{j,z}$: funkcja samego azymutu tej zależności nie opisuje, a odwzorowanie całego wektora tak.
Tłumienie fazy ściska w tym równaniu składowe $y$ i $z$ w sposób deterministyczny, dlatego przy promieniu $0{,}25$--$0{,}41$ wariant AR odzyskuje $53\%$ (M2) i $85\%$ (M1), a wariant FZ $-1\%$ i $73\%$.
Depolaryzacji korekta nie odwraca: próg spadku wypada tam, gdzie promień maleje poniżej $0{,}1$, odzysk wariantów AR i FZ mieści się między $-14\%$ a $+32\%$, wariant AO pogarsza wynik, a odzysk $67\%$, który wariant AR daje w modelu M2 dla wartości dokładnych ($\qdep=0{,}2$, promień $0{,}007$), nie jest osiągalny w pomiarze.
% ŹRÓDŁO: progi.csv, strzaly = 1024 (odzysk_a2, odzysk_b, odzysk_b2, delta_b2_ci_lo, delta_b2_ci_hi, promien_analityczny); pominięte w tekście: w przebiegu mieszanym decyduje składowa, która pierwsza obniża dokładność (M3: odstrojenie, odzysk AR 99%; M1 i M2: depolaryzacja, odzysk AR 3% i -5%); dla strzaly 8192 i 0 odzysk wariantu AR przy odstrojeniu wynosi 75--99%; krzywe.csv, model c7_f200_q3_b13_d2, depolaryzacja 0,2, wariant b2_afin_wprost: strzaly 0 (dokl 0,3934, odzysk_dokl 0,668, promien_analityczny 0,0072) i strzaly 1024 (dokl 0,153 i zgod 0,4455, czyli stała predykcja, tak samo w kontroli z przestawionymi pomiarami w pliku kontrole.csv).
Kalibracja wariantu AR na odłożonej połowie części testowej zamiast na próbkach uczących zmienia jego dokładność przy odstrojeniu i tłumieniu fazy o $-0{,}010$ do $+0{,}003$, więc zagrożenie~(\ref{itm:cls:ZT5}) nie wpływa na te wyniki.
% ŹRÓDŁO: czulosc_kalibracji.csv, wariant b2_afin_wprost, strzaly 1024, poziomy progowe z tabeli: zbior_kalibracyjny uczacy (n_kal 400) oraz testowy_krzyzowo (n_kal 400); różnice policzone skryptem korekta_odczytu.py (ten sam skrypt: kalibracja na 100 zamiast 400 próbek uczących zmienia dokładność najwyżej o 0,009). Pominięte w tekście: wariant AO jest w modelu M2 niestabilny między powtórzeniami kalibracji (krzywe.csv, dokl_sd_powt do 0,0647 wobec 0,0078 dla wariantu AR).
Po przestawieniu pomiarów między próbkami testowymi zgodność predykcji wariantu FO przewyższa zgodność najlepszej stałej predykcji o ponad $0{,}02$ w $12$, $31$ i $1$ ze $117$ ustawień modeli M2, M1 i M3, a pozostałych wariantów w żadnym, więc jego pozornie najlepszy odzysk w modelu M1 ($90\%$ przy odstrojeniu) nie jest wynikiem korekty pomiaru.
% ŹRÓDŁO: kontrole.csv (zgod_tasowany_pomiar wariantu minus większa z wartości zgod_tasowany_pomiar wariantu bez_korekty i zgod_stala_predykcja; największe nadwyżki 0,209; 0,597 i 0,106; zliczenia skryptem korekta_odczytu.py, zgodne z tabelą 6b w README eksperymentu); progi.csv, odzysk_a = 0,90 (c7_f200_q6_b15, odstrojenie).

\paragraph{Niedopasowanie kalibracji.}
Wariant AR skalibrowany przy odstrojeniu do dwóch razy słabszym od rzeczywistego zachowuje zysk, a skalibrowany przy silniejszym traci go i szkodzi tam, gdzie szum jest słaby: w modelu M2 przy $\cohangle=0{,}1$ (bez korekty $0{,}225$) kalibracja przy $0{,}05$ i $0{,}1$ daje $0{,}422$, a przy $0{,}2$ daje $0{,}282$, natomiast przy $\cohangle=0{,}02$ (bez korekty $0{,}440$) kalibracja przy $0{,}2$ daje $0{,}167$.
% ŹRÓDŁO: eksperymenty/krzywe-degradacji/klasyfikator/wyniki/niedopasowanie.csv, skladowa odstrojenie, strzaly 1024, warianty b2_afin_wprost i bez_korekty, kolumny sila_kalibracji, sila_testu, dokl (5 powtórzeń); pominięte w tekście: w modelu c7_f200_q6_b15 zysk znika już na sąsiednim poziomie siatki (przy sile 0,2: bez korekty 0,142, kalibracja zgodna 0,321, kalibracja przy 0,1 i 0,3 daje 0,166 i 0,195); model c7_f200_q3_b13_d2 przy sile 0,1: kalibracja przy 0,03 daje 0,371, przy 0,07 daje 0,427; model c6_f300_q5_b3 przy sile 0,3 (bez korekty 0,328): kalibracja przy 0,1--0,3 daje 0,453--0,459, przy 0,5 daje 0,412.
Nieaktualna korekta bywa więc gorsza od braku korekty, co odpowiada warunkowi~\eqref{eq:pdu:criterion}.

\paragraph{Wynik.}
Odpowiedź na drugą część pytania rozdziału (podrozdział~\ref{sec:cls:scope}) jest warunkowa, a w odniesieniu do urządzenia negatywna: na zapisanych pomiarach korekta odczytu nie zmieniła dokładności w żadnym z dziesięciu modeli, bo w modelach płytkich urządzenie zgadzało się z symulatorem, a w jedynym głębszym modelu z danymi kalibracyjnymi przyrost $0{,}04$--$0{,}05$ nie odróżnia się od zera.
W symulacji korekta przywraca dokładność tylko wtedy, gdy szum zmienia kierunek wektora Blocha systematycznie, a promień pozostaje mierzalny, czyli przy odstrojeniu i częściowo przy tłumieniu fazy; nie odwraca utraty czystości stanu przy skończonej liczbie pomiarów, a skalibrowana przy innym szumie niż rzeczywisty może szkodzić.
Skuteczna jest przy tym regresja całego wektora Blocha, a nie funkcja błędu fazy: wariant FZ, najbliższy funkcji PDU spośród wariantów niekorzystających z informacji o odpowiedzi, odzyskuje przy odstrojeniu $11$--$32\%$ utraconej dokładności, więc funkcja PDU w postaci z rozdziału~\ref{chapter:pdu} nie przenosi się wprost na klasyfikator z odczytem fazowym; przenosi się zasada klasycznej korekty kalibrowanej z pomiaru.
Brak poprawy na urządzeniu jest zgodny z pracami, w których sama mitygacja błędu odczytu nie dała spójnej poprawy klasyfikacji~\cite{peters_machine_2021,singh_benchmarking_2026}.
% ŹRÓDŁO: eksperymenty/przeglad-literatury/sekcja-3.7-klasyfikacja.md, tabela E (prace czytane w pełnym tekście: peters_machine_2021, dod. E; singh_benchmarking_2026, tab. 4).
Poza zagrożeniami (\ref{itm:cls:ZT1})--(\ref{itm:cls:ZT5}) korekta ma dwa ograniczenia własne: każdy wariant wymaga klasycznej symulacji obwodów kalibracyjnych, a wynik z urządzenia opiera się na jednej realizacji przebiegu kalibracyjnego, wykonanego po przebiegu testowym.
% Pominięte w tekście (zagrożenie ZT4 opisuje to ogólnie): eksperymenty/krzywe-degradacji/klasyfikator/wyniki/porownanie_ze_sprzetem.csv, model c7_f200_q3_b13_d2: promień zmierzony na urządzeniu odpowiada krotności 11,17 poziomu ibm_fez (krotnosc_rownowazna), przy której symulowana zgodność z predykcją bez szumu wynosi 0,933 (zgod_sym_przy_krotnosci_rownowaznej_1024) wobec 0,71 na urządzeniu (zgod_sprzet_z_sym_bez_szumu).

\subsection{Trening w szumie a korekta po fakcie}\label{subsec:cls:results:training}
\agentinfo{
Napisane 2026-11-27 na prośbę autora z wyników eksperymentu E2 w dwóch konfiguracjach: X5 (konfiguracja płytka) \agentref{eksperymenty/trening-w-szumie/README.md} i X5b (konfiguracja głęboka) \agentref{eksperymenty/trening-w-szumie/gleboka/README.md}; wszystkie trzy etapy obu eksperymentów (opis, implementacja, wynik) zaakceptował autor.
Liczby pochodzą z plików wyniki/komorki.csv, porownania.csv, dokladnosc.csv i mlp.csv obu katalogów; zaokrąglenia jak w 8.4.1.
Nazwa w tekście: eksperyment porównawczy (opis metodyki w akapicie sekcji 8.3); konfiguracja płytka = konfiguracja przeglądu szumu (K = 3, n = 2, L = 7), konfiguracja głęboka = model M2 uczony od nowa.
Odpowiedniki nazw z plików: poziom T/E 1, 2, 3 = poziomy stopniowane 0,01; 0,05; 0,10 wszystkich trzech składowych; E4 = poziom ibm_fez; wariant 1 = bez działań, wariant 2 = korekta po fakcie, wariant 3 = trening w szumie, wariant 4 = oba; a2 = FZ, b2 = AR, a3 = FZB (nowy wariant X5b: FZ uzupełniony o funkcję błędu kąta biegunowego z argumentem zmierzonym, tylko analiza uzupełniająca).
Rodzina główna była ustalona w opisach przed obliczeniami: dla konfiguracji płytkiej poziom 0,10, dla głębokiej 0,05; wyraźne efekty przy 0,10 w konfiguracji głębokiej są analizą uzupełniającą i tak je opisano; nie przenosić ich do rodziny głównej.
Etykieta dodana: tab:cls:training.
Objętość w kompilacji bez notatek (2026-11-27): podrozdział 8.4.3 około 3 strony (str. 106--108, z tabelą około 0,5 strony), akapit Eksperyment porównawczy w 8.3 około 0,4 strony; rozdział 8 ma teraz 20 stron (str. 91--110) przy budżecie 16, a podrozdział 8.4.4 i sekcje 8.5--8.6 są puste; skracanie czeka na sygnał autora. Kandydaci do cięcia w 8.4.3: wiersze składowych w tabeli (0,15 strony), akapit Punkt odniesienia klasyczny (przenieść jedno zdanie do 8.4.1; 0,1), zdania o ziarnach w akapicie Rodzina główna (0,15), zdanie o koszcie treningu w szumie (0,05).
Pominięte dla budżetu stron (dostępne w plikach): kalibracja na części walidacyjnej (w konfiguracji głębokiej zysk AR przy 0,05 znika: +0,001), pełne macierze niedopasowania, wariant FZB poza zdaniem o odstrojeniu, kontrole negatywne poza jednym zdaniem, testy McNemara poza dwoma ziarnami.
}

Eksperyment porównawczy (podrozdział~\ref{sec:cls:methodology}) zestawia korektę odczytu z dwóch poprzednich podrozdziałów z drugą strategią obrony przed szumem: uczeniem sieci $\Theta_{\mathbf{w}}$ z modelem szumu w obwodzie, tak by sama nauczyła się go kompensować.
Obie strategie porównano na dwóch konfiguracjach uczonych od nowa na pięciu ziarnach: płytkiej, czyli konfiguracji przeglądu szumu ($K=3$, $d=200$, $n=2$, $L=7$, wariant Pauliego), i głębokiej, czyli modelu M2 ($K=7$, $d=200$, $n=3$, $L=13$, wariant uniwersalny, dwie współrzędne z kubitu).
Przy szumie testu o danym poziomie porównywano cztery strategie: bez działań (model uczony bez szumu), korektę po fakcie (ten sam model z korektą odczytu FZ albo AR), trening w szumie (model uczony na poziomie szumu testu) oraz obie naraz.
Rodzinę główną pięciu porównań ustalono przed obliczeniami dla $\nshots=1024$: na poziomie $0{,}10$ w konfiguracji płytkiej i $0{,}05$ w głębokiej, bo według danych modelu M2 z podrozdziału~\ref{subsec:cls:results:noise} rzut wektora Blocha na płaszczyznę $xy$ maleje przy $0{,}10$ do około $0{,}05$, czyli do wielkości porównywalnej z niepewnością statystyczną składowej ($1/\sqrt{1024}\approx0{,}03$).
% ŹRÓDŁO: eksperymenty/trening-w-szumie/OPIS.md (punkt 9) i gleboka/OPIS.md (punkt 9, uzasadnienie wyboru poziomu 0,05: summary.json odtworzeń modelu M2, m_mean 0,149 i 0,047 przy poziomach 0,05 i 0,10).

\paragraph{Spadek dokładności bez działań.}
W konfiguracji płytkiej poziomy $0{,}01$ i $0{,}05$ zmieniają dokładność modelu uczonego bez szumu najwyżej o $0{,}001$, a poziom $0{,}10$ obniża ją średnio o $0{,}048$ $[-0{,}156;\,+0{,}060]$, przy czym spadek pochodzi niemal w całości z jednego ziarna, w którym błąd azymutu wynosi $45^{\circ}$ wobec $13^{\circ}$--$18^{\circ}$ w pozostałych, a dokładność $0{,}556$ wobec $0{,}781$ bez szumu (wartości dokładne).
W konfiguracji głębokiej poziom $0{,}01$ i poziom ibm\_fez nie zmieniają dokładności (najwyżej o $0{,}001$), poziom $0{,}05$ obniża ją o $0{,}053$ $[-0{,}113;\,+0{,}007]$ w czterech z pięciu ziaren, a poziom $0{,}10$ o $0{,}202$ $[-0{,}249;\,-0{,}155]$ we wszystkich, do $0{,}163$ wobec $1/7\approx 0{,}143$ dla wyboru losowego.
Potwierdza to na pięciu ziarnach i pełnej części testowej wynik pojedynczego modelu M2 z podrozdziału~\ref{subsec:cls:results:noise}.
Dokładności bezwzględne są niższe niż w tabeli~\ref{tab:cls:hw-results} ($0{,}364$ wobec $0{,}439$), bo punkt kontrolny wybiera się tu na części walidacyjnej, a część testowa jest inna~(\ref{itm:cls:ZT1}).
% ŹRÓDŁO: eksperymenty/trening-w-szumie/wyniki/porownania.csv, rodzina warunek_wstepny_H0, tryb strzaly (E1 -0,0002; E2 -0,0012; E3 -0,0482 [-0,1563; +0,0599]) i dokladnosc.csv (T0, tryb dokladny, ziarno 43: E0 0,781, E3 0,556, faza_st 44,9; pozostałe ziarna przy E3: faza_st 13,3--17,5); gleboka/wyniki/porownania.csv, warunek_wstepny_H0, strzaly (E1 -0,0002; E2 -0,0526 [-0,1126; +0,0074], 1/4 ziaren; E3 -0,2024 [-0,2493; -0,1554], 0/5; E4 -0,0010) i komorki.csv (T0, strzaly, E3 0,1626; T0, dokladny, E0 0,3638).

\paragraph{Rodzina główna.}
W żadnej konfiguracji żadne z pięciu porównań nie ma przedziału ufności bez zera, a $p$-wartości po poprawce Holma wynoszą $1$ w konfiguracji płytkiej i od $0{,}45$ do $0{,}62$ w głębokiej (tabela~\ref{tab:cls:training}).
Względem braku działań trening w szumie daje $+0{,}046$ $[-0{,}070;\,+0{,}162]$ i $+0{,}059$ $[-0{,}025;\,+0{,}143]$, wariant AR $+0{,}049$ $[-0{,}059;\,+0{,}156]$ i $+0{,}026$ $[-0{,}025;\,+0{,}076]$, a wariant FZ $+0{,}031$ $[-0{,}045;\,+0{,}107]$ i $-0{,}003$ $[-0{,}021;\,+0{,}014]$ (kolejno konfiguracja płytka i głęboka); trening w szumie różni się od wariantu AR o $-0{,}003$ $[-0{,}047;\,+0{,}041]$ i $+0{,}033$ $[-0{,}008;\,+0{,}075]$.
Na poziomach rodziny głównej spadek bez działań jest mały i niejednolity, a obie strategie zależą od ziarna.
W konfiguracji płytkiej w ziarnie ze spadkiem każda strategia istotnie poprawia wynik, ale model uczony w szumie jest w innym ziarnie gorszy od modelu bez działań (wartości dokładne: $0{,}740$ wobec $0{,}783$; dokładny test McNemara w obrębie ziarna: $69$ próbek zyskało, $112$ straciło, $p=0{,}033$ po poprawce Holma).
W konfiguracji głębokiej jedno ziarno nie traci dokładności przy $0{,}05$, a wariant AR obniża ją w nim z $0{,}427$ do $0{,}398$ ($17$ zysków, $46$ strat, $p=0{,}007$).
% ŹRÓDŁO: eksperymenty/trening-w-szumie/wyniki/porownania.csv i gleboka/wyniki/porownania.csv, rodzina glowna (roznica, ci95_lo, ci95_hi, p_t_holm_w_rodzinie_glownej; płytka: 3-2a2 +0,015 [-0,036; +0,066]; głęboka: 3-2a2 +0,062 [-0,016; +0,140]) oraz rodzina mcnemar_w_ziarnie, tryb dokładny (płytka: E3, ziarno 43: zyski wszystkich trzech strategii 304/78, 260/40, 291/63; ziarno 45, wariant 3 - wariant 1: 69/112, p Holma 0,033; dokladnosc.csv, ziarno 45, dokladny, E3: T3 0,740, T0 0,783; głęboka: E2, ziarno 45, wariant 2 b2 - wariant 1: 17/46, p Holma 0,0071; dokladnosc.csv, ziarno 45, dokladny, E2: bez działań 0,427, b2 0,398; przy 1024 pomiarach 0,421 i 0,387).

\begin{table}[tb]
\centering
\caption[Trening w szumie a korekta po fakcie: dokładność na pięciu ziarnach]{Trening w szumie a korekta po fakcie w eksperymencie porównawczym: dokładność na części testowej ($1000$ próbek na ziarno), średnia $\pm$ odchylenie standardowe po pięciu ziarnach, $\nshots=1024$, w wierszach oznaczonych (d) wartości dokładne. Szum testu: wszystkie trzy składowe na poziomie stopniowanym albo jedna składowa o sile $0{,}10$. Bez działań: model uczony bez szumu; FZ i AR: ten sam model z korektą odczytu kalibrowaną na $3000$ próbkach uczących; trening w szumie: model uczony na poziomie szumu testu. Konfiguracja płytka: $K=3$, $n=2$, $L=7$; głęboka: model M2, $K=7$, $n=3$, $L=13$; odstr.: odstrojenie, depol.: depolaryzacja. Wiersze rodziny głównej oznaczono gwiazdką.}
\label{tab:cls:training}
\footnotesize
\setlength{\tabcolsep}{3pt}
\begin{tabular}{@{}llcccc@{}}
\toprule
Konf. & Szum testu & Bez działań & FZ & AR & Trening w szumie \\
\midrule
płytka & brak & $0{,}785\pm0{,}012$ & -- & -- & -- \\
 & $0{,}01$ & $0{,}785\pm0{,}012$ & $0{,}785\pm0{,}012$ & $0{,}785\pm0{,}012$ & $0{,}781\pm0{,}006$ \\
 & $0{,}05$ & $0{,}783\pm0{,}010$ & $0{,}783\pm0{,}010$ & $0{,}785\pm0{,}012$ & $0{,}786\pm0{,}021$ \\
 & $0{,}10$\,$^{\ast}$ & $0{,}737\pm0{,}089$ & $0{,}768\pm0{,}029$ & $0{,}786\pm0{,}012$ & $0{,}783\pm0{,}024$ \\
 & odstr. $0{,}10$ & $0{,}770\pm0{,}026$ & $0{,}782\pm0{,}007$ & $0{,}785\pm0{,}013$ & -- \\
 & depol. $0{,}10$ & $0{,}784\pm0{,}012$ & $0{,}784\pm0{,}012$ & $0{,}786\pm0{,}012$ & -- \\
\addlinespace
głęboka & brak & $0{,}365\pm0{,}039$ & -- & -- & -- \\
(M2) & $0{,}01$ & $0{,}365\pm0{,}039$ & $0{,}364\pm0{,}038$ & $0{,}364\pm0{,}038$ & $0{,}374\pm0{,}032$ \\
 & $0{,}05$\,$^{\ast}$ & $0{,}312\pm0{,}071$ & $0{,}309\pm0{,}060$ & $0{,}338\pm0{,}043$ & $0{,}371\pm0{,}041$ \\
 & $0{,}10$ & $0{,}163\pm0{,}018$ & $0{,}177\pm0{,}028$ & $0{,}171\pm0{,}010$ & $0{,}310\pm0{,}043$ \\
 & $0{,}10$ (d) & $0{,}187\pm0{,}063$ & $0{,}246\pm0{,}098$ & $0{,}334\pm0{,}043$ & $0{,}367\pm0{,}055$ \\
 & odstr. $0{,}10$ & $0{,}187\pm0{,}056$ & $0{,}229\pm0{,}043$ & $0{,}360\pm0{,}043$ & -- \\
 & depol. $0{,}10$ & $0{,}223\pm0{,}037$ & $0{,}239\pm0{,}040$ & $0{,}217\pm0{,}039$ & -- \\
 & depol. $0{,}10$ (d) & $0{,}346\pm0{,}040$ & $0{,}347\pm0{,}038$ & $0{,}357\pm0{,}034$ & -- \\
\bottomrule
\end{tabular}
\end{table}
% ŹRÓDŁO tabeli: eksperymenty/trening-w-szumie/wyniki/komorki.csv (konfiguracja płytka) i eksperymenty/trening-w-szumie/gleboka/wyniki/komorki.csv (głęboka): tryb strzaly albo dokladny, kalibracja uczaca, T = 0 (bez działań, FZ = korekta a2, AR = korekta b2) albo T = poziom testu (trening w szumie, korekta brak), E = E0, E1, E2, E3, det3, dep3; kolumny dokl_srednia, dokl_sd.

\paragraph{Składowe szumu.}
Każdą składową o sile $0{,}10$ podawano też osobno, bez nowego uczenia (analiza uzupełniająca, $p$ bez poprawki).
W obu konfiguracjach szkodzi przede wszystkim odstrojenie, czyli błąd koherentny, a wariant AR usuwa jego skutek całkowicie: w konfiguracji głębokiej dokładność rośnie z $0{,}187$ do $0{,}360$ (o $0{,}174$ $[+0{,}100;\,+0{,}247]$, we wszystkich ziarnach), przy błędzie azymutu $34^{\circ}$ i kąta biegunowego $41^{\circ}$ oraz prawie niezmienionym promieniu ($0{,}79$).
Wariant FZ odzyskuje tu tylko $0{,}04$, bo nie zmienia kąta biegunowego, a odstrojenie jest obrotem wokół osi $x$ (równanie~\eqref{eq:cls:noise-frame}); dopisany do niego wariant FZB, w którym funkcję błędu z argumentem zmierzonym wyznacza się także dla kąta biegunowego, odzyskuje $0{,}12$.
Depolaryzacja skraca wektor w konfiguracji płytkiej do promienia $0{,}31$, prawie bez zmiany decyzji, a w głębokiej do $0{,}06$; tam decyzję zmienia głównie szum próbkowania: dokładność wynosi $0{,}346$ dla wartości dokładnych, $0{,}308$ przy $8192$ i $0{,}223$ przy $1024$ pomiarach, a żaden wariant korekty jej nie przywraca.
Tłumienie fazy o sile $0{,}10$ obniża dokładność najwyżej o $0{,}007$.
Odpowiada to wynikom eksperymentu symulacyjnego (tabela~\ref{tab:cls:pdu-recovery}), tym razem na pięciu ziarnach.
% ŹRÓDŁO: gleboka/wyniki/porownania.csv, rodzina kanaly_osobno, strzaly, det3 (b2: +0,1736 [+0,1001; +0,2471], 5/0; a2: +0,0424; a3: +0,1241), dep3 (b2: -0,0062); gleboka/wyniki/tabela.md, tabela 7 (det3: błąd azymutu 33,6, kąta biegunowego 41,2, r 0,787; dep3: r 0,062); komorki.csv, dep3 bez korekty: dokladny 0,3458, strzaly8192 0,3078, strzaly 0,2230; deph3: 0,358 (strzaly, głęboka) wobec 0,365 przy E0; konfiguracja płytka: README.md, tabela mechanizmu (depolaryzacja poziom 3: r 0,306; dokl. 0,785) i komorki.csv (det3, deph3, dep3).

\paragraph{Silny szum w konfiguracji głębokiej.}
Przy wszystkich trzech składowych o sile $0{,}10$ i $1024$ pomiarach pomaga tylko trening w szumie: $0{,}310$ wobec $0{,}163$ bez działań (o $0{,}147$ $[+0{,}095;\,+0{,}199]$, we wszystkich ziarnach), podczas gdy obie korekty zostają przy wyborze losowym ($0{,}177$ i $0{,}171$; analiza uzupełniająca).
Informacja o klasie nie znika jednak ze stanów: dla wartości dokładnych wariant AR podnosi dokładność do $0{,}334$, a przy $8192$ pomiarach do $0{,}244$.
Ten zysk wymaga prawdziwego powiązania odczytów, bo korekta AR dopasowana do par o przestawionym przyporządkowaniu daje wybór losowy ($0{,}151$), natomiast część zysku wariantu FZ ($0{,}187\to0{,}246$) daje także korekta dopasowana do par przestawionych ($0{,}221$).
Rzut wektora na płaszczyznę $xy$ wynosi wtedy około $0{,}045$, więc według oszacowania~\eqref{eq:cls:azimuth-noise} niepewność azymutu przy $1024$ pomiarach sięga około $0{,}7$~rad, a takiego rozrzutu korekta kalibrowana po fakcie nie odwraca.
Model uczony w szumie zachowuje $0{,}310$, choć jego wektory są niewiele dłuższe ($0{,}082$ wobec $0{,}049$); przyczyny tej różnicy eksperyment nie rozstrzyga.
Korekta AR dodana do treningu w szumie obniża wtedy wynik do $0{,}262$.
Trening w silnym szumie ma też koszt: model uczony na poziomie $0{,}10$ ma bez szumu dokładność $0{,}278$ wobec $0{,}365$, a uczony na poziomie $0{,}05$ zachowuje $0{,}367$ bez szumu i $0{,}371$ przy $0{,}05$, lecz przy $0{,}10$ spada do $0{,}276$.
% ŹRÓDŁO: gleboka/wyniki/tabela.md, tabela 8 (kontrole negatywne, model T0, tryb dokładny, E3: bez korekty 0,187; a2 po korekcie 0,246, tasowane pary 0,221; b2 po korekcie 0,334, tasowane pary 0,151; tasowany odczyt 0,140--0,149) i kontrole.csv; gleboka/wyniki/porownania.csv, rodzina uzupelniajace, strzaly, uczaca, E3, wariant 3 - wariant 1 (+0,1468 [+0,0950; +0,1986], 5/0); komorki.csv: E3 dokladny b2 0,3340; strzaly8192 b2 0,2442; strzaly T3 b2 0,2620; tryb strzaly: T3 przy E0 0,27765, T0 przy E0 0,36504, T2 przy E0 0,36687, przy E2 0,3713 i przy E3 0,2763; dokladnosc.csv, średnie m i r modelu T0 przy E3 (0,045; 0,049) i T3 przy E3 (r 0,082), tryb dokladny; 1/(0,045 * 32) = 0,69 rad.

\paragraph{Punkt odniesienia klasyczny.}
Sieć MLP o tych samych warstwach ukrytych, uczona i wybierana na tych samych podziałach, osiąga $0{,}793\pm0{,}016$ w konfiguracji płytkiej i $0{,}496\pm0{,}018$ w głębokiej, wobec $0{,}785\pm0{,}012$ i $0{,}364\pm0{,}040$ klasyfikatora bez szumu.
W konfiguracji płytkiej jest lepsza w czterech z pięciu ziaren (o $0{,}008$, $p=0{,}44$), w głębokiej we wszystkich (o $0{,}132$ $[+0{,}080;\,+0{,}184]$, $p=0{,}002$).
% ŹRÓDŁO: wyniki/mlp.csv obu katalogów (dokl_testowa, klasyfikator_kwantowy_T0_E0_dokladny); test t sparowany po ziarnach policzony dodatkowo (głęboka: różnice 0,126; 0,184; 0,159; 0,113; 0,076, p = 0,0021); płytka: README.md (różnica +0,008, p = 0,44).

\paragraph{Wynik.}
Testy zaplanowane z góry nie rozstrzygają, która strategia lepiej chroni klasyfikator: na poziomach rodziny głównej szum obniża dokładność słabo i niejednolicie, a obie strategie poprawiają wynik w części ziaren i pogarszają w innych.
Analizy uzupełniające pokazują podział zadań zgodny z podrozdziałem~\ref{subsec:cls:results:pdu}: korekta po fakcie w wariancie AR usuwa błąd koherentny, dopóki rzut wektora Blocha jest wyraźnie większy od niepewności statystycznej, a gdy depolaryzacja w głębokim obwodzie skraca wektor do wielkości tej niepewności, pomaga tylko trening w szumie, kosztem dokładności przy słabszym szumie.
Wariant FZ, najbliższy funkcji PDU, jest najsłabszy, bo koryguje sam azymut.
Wnioski dotyczą symulatora z uproszczonym modelem szumu~(\ref{itm:cls:ZT4}), dwóch konfiguracji uczonych na jednym zbiorze każda i pięciu ziaren; architektury i hiperparametry pochodzą ze strojenia na zbiorze testowym~(\ref{itm:cls:ZT1}), a wyraźne efekty przy poziomie $0{,}10$ w konfiguracji głębokiej są analizą uzupełniającą, bo rodzinę główną tej konfiguracji ustalono na poziomie $0{,}05$.

\subsection{Klasyfikacja obrazów po redukcji błędów}\label{subsec:cls:results:images}
\agentinfo{Eksperyment E4. Pisać po jego wykonaniu. Zwrócić uwagę na rozłączność danych użytych do treningu sieci pix2pix i do testu klasyfikatora.}

\section{Dyskusja}\label{sec:cls:discussion}
\agentinfo{
Odpowiedź na PB5 oraz zagrożenia trafności.
Co korekta klasyczna usuwa (składowa koherentna, systematyczne przesunięcie fazy), a czego nie (utrata czystości stanu).
}

\section{Podsumowanie}\label{sec:cls:summary}
